% ---------------------------------------------------------------------------
% Multi-agent reinforcement learning algorithms (marl_algorithms)
% ---------------------------------------------------------------------------
\begingroup
% Line breaking for paragraphs with long, unbreakable identifiers (local to
% this chapter): allow looser lines instead of overfull ones.
\setlength{\emergencystretch}{3em}\tolerance=1000
\chapter{Multi-Agent Reinforcement Learning Algorithms}\label{chap:marl}

The package \pkg{marl\_algorithms}, new in version 1.2, contains reference
implementations of seven classical algorithms for cooperative multi-agent
reinforcement learning: independent and multi-agent proximal policy
optimisation (IPPO, MAPPO), multi-agent deep deterministic policy gradient and
its twin-critic variant (MADDPG, MATD3), and independent Q-learning, value
decomposition networks and QMIX (IQL, VDN, QMIX). All seven are written
against one small core and train directly on the environments of
\pkg{env\_lib}: on the native batched vector environments where these exist,
with the per-agent rewards that every environment reports, and with evaluation
against the classical baseline controllers of
Section~\ref{sec:workflow-baselines}.

The code is meant to be read next to the papers. Every algorithm module cites
the method it implements, every class docstring states the loss it minimises,
and every simplification or deviation from the original method is stated in
the code and in this chapter. Tuned \emph{presets} train each algorithm on one
or more demonstration environments in at most about two minutes on a single
CPU core.
The package requires PyTorch (the \code{torch} extra,
Section~\ref{sec:install-extras}); \code{import marl\_algorithms} itself loads
nothing but the version number, and the first use of an algorithm imports
PyTorch.

% ---------------------------------------------------------------------------
\section{Overview}\label{sec:marl-overview}

Three parts of the package connect the algorithms to \pkg{env\_lib}:

\begin{enumerate}
  \item \emph{A per-agent view of the joint spaces.} The environments expose
    the whole team through joint spaces: observations stacked as
    \code{(n\_agents, obs\_dim)} and one joint action
    (Section~\ref{sec:envlib-joint}). \code{MultiAgentSpec} reads the number of
    agents, the size of one agent's observation and the kind and size of one
    agent's action from these spaces, and converts observations to per-agent
    arrays and per-agent actions back to the joint action
    (Section~\ref{sec:marl-spec}).
  \item \emph{Training on batched copies.} \code{make\_vector\_env} creates
    \code{num\_envs} copies of an environment with \code{env\_lib.make\_vec},
    which uses the native batched implementation when there is one
    (Section~\ref{sec:workflow-vector}), with the \code{"same\_step"}
    autoreset mode, so that every transition is a real one and the final
    observation of a truncated episode is available for bootstrapping. A
    \code{VectorRunner} turns every vector step into a \code{Transition} of
    per-agent arrays, including the per-agent rewards of
    \code{info["agent\_rewards"]} (Section~\ref{sec:marl-runner}).
  \item \emph{Evaluation with the \pkg{env\_lib} tools.} The method
    \code{policy()} of a trained algorithm returns a function from environment
    observations to environment actions that works for single and vector
    environments. The evaluation function \code{env\_lib.evaluate} therefore
    compares the trained policy with random actions and with the baseline
    controller on the same seeded episodes, and the recording function
    \code{record\_episode} of \pkg{env\_lib.utils} turns an episode of it
    into an animation (Section~\ref{sec:marl-quickstart}).
\end{enumerate}

Table~\ref{tab:marl-algorithms} lists the algorithms as registered in
\pkg{marl\_algorithms.registry}; the function \code{list\_algorithms()}
returns these records and the command \code{marl-train list} prints them.
Section~\ref{sec:marl-references} gives the full references cited in the
source code.

\begin{table}[htbp]
  \centering\small
  \caption{Registered algorithms. ``Family'' and ``Actions'' are the
    \code{family} and \code{action\_kinds} fields of \code{AlgorithmInfo}.}
  \label{tab:marl-algorithms}
  \begin{tabularx}{\textwidth}{@{}>{\raggedright\arraybackslash}p{1.35cm}>{\raggedright\arraybackslash}p{2.1cm}>{\raggedright\arraybackslash}p{1.8cm}>{\raggedright\arraybackslash}X>{\raggedright\arraybackslash}p{2.2cm}@{}}
    \toprule
    Name & Family & Actions & Key idea & Reference \\
    \midrule
    \code{ippo} & on-policy & continuous, discrete &
      PPO for every agent; decentralised critic $V(o_i)$ trained on the
      agent's own reward & de Witt et al.\ (2020) \\
    \code{mappo} & on-policy & continuous, discrete &
      decentralised PPO actors; centralised critic $V(s, i)$ on the global
      state, trained on the team reward & Yu et al.\ (2022) \\
    \code{maddpg} & off-policy actor-critic & continuous &
      deterministic actors $\mu_i(o_i)$; centralised critics
      $Q_i(s, a_1, \dots, a_n)$ on the global state and the joint action &
      Lowe et al.\ (2017) \\
    \code{matd3} & off-policy actor-critic & continuous &
      MADDPG with twin critics, target policy smoothing and delayed actor
      updates & Ackermann et al.\ (2019) \\
    \code{iql} & value-based & discrete &
      one DQN utility $Q_i(o_i, a_i)$ per agent, each trained on its own
      reward & Tan (1993) \\
    \code{vdn} & value-based & discrete &
      team value $\sum_i Q_i(o_i, a_i)$, trained end to end on the team
      reward & Sunehag et al.\ (2018) \\
    \code{qmix} & value-based & discrete &
      team value as a monotonic, state-conditioned mixture of the utilities
      (hypernetworks) & Rashid et al.\ (2018) \\
    \bottomrule
  \end{tabularx}
\end{table}

Every algorithm works on every environment whose action kind it supports.
Table~\ref{tab:marl-envs} shows the per-agent structure that
\code{MultiAgentSpec} reads from the registered environments; every
environment was trained briefly with an algorithm of each family that applies
to it, and an algorithm given an environment of the wrong action kind raises
\code{TypeError}.
Environments whose observation is not stacked per agent, the Kuramoto
environments, are treated as a single agent that observes and acts on the whole
network; with one agent IPPO and MAPPO reduce to PPO, MADDPG and MATD3 to DDPG
and TD3, and IQL to (double) DQN.

\begin{table}[htbp]
  \centering\small
  \caption{Per-agent structure of the registered environments as read by
    \code{MultiAgentSpec} (default parameters) and the algorithms that apply.
    $n$ is the number of agents and $d$ the size of one agent's observation.}
  \label{tab:marl-envs}
  \begin{tabularx}{\textwidth}{@{}>{\raggedright\arraybackslash}p{4.3cm}rr>{\raggedright\arraybackslash}p{3.1cm}>{\raggedright\arraybackslash}X@{}}
    \toprule
    Environment & $n$ & $d$ & Action per agent & Algorithms \\
    \midrule
    \envid{Consensus-v0}, \envid{Formation-v0} & 8 & 16 & continuous, 2-d & IPPO, MAPPO, MADDPG, MATD3 \\
    \envid{PowerGrid-v0} & 16 & 10 & continuous, 1-d & IPPO, MAPPO, MADDPG, MATD3 \\
    \envid{Platoon-v0} & 8 & 13 & continuous, 1-d & IPPO, MAPPO, MADDPG, MATD3 \\
    \envid{AJLATT-v0} & 4 & 24 & continuous, 2-d & IPPO, MAPPO, MADDPG, MATD3 \\
    \envid{Pistonball-v0} & 20 & 7 & continuous, 1-d; 3 choices with \code{continuous=False} & all seven, by action mode \\
    \envid{LineMsg-v0} (per-agent bits) & 10 & 3 & 2 choices & IPPO, MAPPO, IQL, VDN, QMIX \\
    \envid{WirelessComm-v1} (\envid{-v0}) & 16 (36) & 18 & 5 choices & IPPO, MAPPO, IQL, VDN, QMIX \\
    \envid{KuramotoOscillator-v0} (other Kuramoto ids: other sizes) & 1 & 75 & continuous, 55-d & IPPO, MAPPO, MADDPG, MATD3 \\
    \bottomrule
  \end{tabularx}
\end{table}

% ---------------------------------------------------------------------------
\section{Quick start}\label{sec:marl-quickstart}

\subsection{Training and evaluating from Python}

\code{marl\_algorithms.train} trains a registered algorithm on batched copies
of an environment and returns the trained algorithm together with its
\code{TrainingLog}. The run below is deliberately short (2\,000 environment
steps, a few seconds); the tuned MAPPO preset for \envid{PowerGrid-v0} trains
for 250\,000 steps (Section~\ref{sec:marl-results}).

\begin{pycode}
import env_lib
from marl_algorithms import train

# A short run; the tuned MAPPO preset for PowerGrid-v0 trains for 250_000 steps.
algo, log = train("mappo", "PowerGrid-v0", total_steps=2_000, num_envs=8,
                  seed=0)
print(algo)
print(len(log.episodes), "training episodes, mean return",
      round(log.mean_return(), 1))

envs = env_lib.make_vec("PowerGrid-v0", 16)
for name, policy in [("random", None), ("mappo", algo.policy()),
                     ("droop", env_lib.baseline_policy(envs))]:
    result = env_lib.evaluate(envs, policy, n_episodes=32, seed=1)
    print(f"{name:8s}{result.mean_return:10.2f}")
envs.close()
\end{pycode}

\begin{textcode}
MAPPO(16 agents, 10-d observations, continuous 1-d per agent, env_steps=2048)
8 training episodes, mean return -381.0
random     -613.22
mappo       -49.52
droop        -0.71
\end{textcode}

% norun
\begin{pycode}
train(algorithm, env_id, total_steps, *, num_envs=16, seed=0, env_kwargs=None,
      device="cpu", callback=None, **config) -> (Algorithm, TrainingLog)
\end{pycode}
creates the copies with \code{make\_vector\_env}, builds the algorithm with the
seed and the configuration overrides \code{**config} (field names of the
algorithm's configuration class; an unknown name raises \code{TypeError}),
trains it with \code{learn}, whose environment reset uses the same seed, and
closes the copies. On-policy algorithms stop after the first
complete rollout that reaches \code{total\_steps}, so the run above collected
2\,048 steps (two rollouts of 64 steps on 8 copies). The training returns in
the log are those of the exploratory policy; the evaluation uses the
deterministic policy (the mean action), which \code{algo.policy()} returns by
default.

The presets of Section~\ref{sec:marl-results} are looked up and trained with the
functions of \pkg{marl\_algorithms.presets}, which the top-level package also
exports:

\begin{pycode}
from marl_algorithms import get_preset, list_presets, train_preset

print(len(list_presets()), list_presets()[:2])
print(get_preset("qmix", "LineMsg-v0"))

# The preset trains for 25_000 steps; a smoke test overrides the length.
algo, log = train_preset("qmix", "LineMsg-v0", total_steps=2_000,
                         warmup_steps=500)
print(algo.config.lr, algo.config.batch_size, algo.env_steps)
\end{pycode}

\begin{textcode}
17 [('ippo', 'PowerGrid-v0'), ('ippo', 'Platoon-v0')]
{'num_envs': 16, 'total_steps': 25000, 'config': {'batch_size': 128,
'lr': 0.001, 'epsilon_decay_steps': 12500}, 'env_kwargs': {}}
0.001 128 2000
\end{textcode}

% norun
\begin{pycode}
train_preset(algorithm, env_id, *, seed=0, total_steps=None, num_envs=None,
             device="cpu", callback=None, **config) -> (Algorithm, TrainingLog)
\end{pycode}
calls \code{train} with the preset's values; \code{total\_steps},
\code{num\_envs} and configuration overrides replace them. It raises \code{KeyError} when the pair has no preset.
\code{get\_preset} returns a copy of the preset (or \code{None}), so it can be
modified freely.

A trained algorithm drives any environment of the same structure through
\code{policy()}, and \code{save} and \code{Algorithm.load} store and restore
it:

\begin{pycode}
import numpy as np
import env_lib
from env_lib.utils import record_episode
from marl_algorithms import Algorithm, make_vector_env, train

algo, log = train("vdn", "LineMsg-v0", total_steps=2_000, warmup_steps=500)

# Record the greedy policy on a single environment with per-agent actions.
env = env_lib.make("LineMsg-v0", action_space_type="multibinary",
                   render_mode="rgb_array")
frames = record_episode(env, algo.policy(), "renders/vdn_linemsg.gif",
                        seed=1, render_every=2)
env.close()

# Save and restore: the class and the configuration are read from the file.
path = algo.save("runs/vdn_linemsg.pt")
restored = Algorithm.load(path)
envs = make_vector_env("LineMsg-v0", 4)
obs = restored.spec.agent_obs(envs.reset(seed=0)[0])     # (4, 10, 3)
assert np.array_equal(restored.act(obs, deterministic=True),
                      algo.act(obs, deterministic=True))
print(type(restored).__name__, restored.env_steps, len(frames))
envs.close()
\end{pycode}

The file written by \code{save} holds the registry name, the
\code{MultiAgentSpec}, the configuration, the seed and the state (networks,
optimisers, counters and normalisation statistics); \code{Algorithm.load}
selects the class through the registry. The replay buffer of the off-policy
methods is not saved. Checkpoints are read with \code{torch.load(...,
weights\_only=False)} because they contain more than tensors, so load only
files from trusted sources.

\paragraph{Training in stages.} The method \code{learn} of an algorithm
(Section~\ref{sec:marl-base}) trains until the counter \code{algo.env\_steps}, which
persists across calls and through \code{save} and \code{load}, reaches
\code{total\_steps}; \code{total\_steps} is therefore a cumulative target, not
a number of additional steps. To continue a run, call \code{learn} again with
a larger total and pass the previous log to append to it. Every call resets
the copies with its \code{seed}, so episodes still running at the end of a
call are not completed. The callback is
called as \code{callback(algorithm, log)} after every update (on-policy) or
every round of gradient steps (off-policy); returning \code{False} stops
training. The environment copies passed to \code{learn} must use the
\code{"same\_step"} autoreset mode, which \code{make\_vector\_env} sets:

\begin{pycode}
from marl_algorithms import MATD3, make_vector_env

def stop_when_solved(algo, log):
    """Stop once the last 20 training episodes average more than -1600."""
    if log.mean_return(20) > -1600:     # NaN (no episode yet) compares False
        return False

envs = make_vector_env("Consensus-v0", 16)
algo = MATD3(envs, seed=0, reward_scale=0.01, normalize_observations=True)
log = algo.learn(envs, 3_200, seed=0)            # one episode per copy
log = algo.learn(envs, 6_400, seed=1, log=log,   # 3_200 more steps
                 callback=stop_when_solved)
print(algo.env_steps, algo.num_updates, len(log.episodes))
log.to_csv("runs/matd3_consensus.csv")
envs.close()
\end{pycode}

\subsection{The \texorpdfstring{\code{marl-train}}{marl-train} command line}\label{sec:marl-cli}

Installing the package provides the console script \code{marl-train}
(equivalently \code{python -m marl\_algorithms}) with the sub-commands of
Table~\ref{tab:marl-cli}. Values given with \code{-{}-set} and
\code{-{}-env-kwarg} are parsed as Python literals (plain text otherwise) and
override the preset's configuration and environment arguments.

\begin{table}[htbp]
  \centering\small
  \caption{Sub-commands of \code{marl-train} (defaults in brackets).}
  \label{tab:marl-cli}
  \begin{tabularx}{\textwidth}{@{}>{\raggedright\arraybackslash}p{0.36\textwidth}>{\raggedright\arraybackslash}X@{}}
    \toprule
    Command & Purpose \\
    \midrule
    \code{list} & the registered algorithms with family, action kinds, summary
      and reference \\
    \code{presets} & every \code{(algorithm, environment)} pair with a preset,
      with its steps and copies \\
    \code{run ALGORITHM ENV\_ID} & train (with the preset when there is one),
      print progress and a summary, optionally save the algorithm and the
      training episodes, then evaluate random actions, the trained policy and
      the baseline \\
    \quad \code{-{}-steps N}, \code{-{}-num-envs N} & environment steps and
      batched copies [preset; 100\,000 and 16 without a preset] \\
    \quad \code{-{}-seed S} & seed of training [0]; the evaluation uses
      \code{S + 1} \\
    \quad \code{-{}-no-preset} & ignore the preset and use the defaults of the
      configuration class \\
    \quad \code{-{}-set KEY=VALUE ...} & configuration overrides \\
    \quad \code{-{}-env-kwarg KEY=VALUE ...} & environment arguments \\
    \quad \code{-{}-eval-episodes N} & evaluation episodes [64]; 0 skips the
      evaluation \\
    \quad \code{-{}-reports N} & progress lines during training [10] \\
    \quad \code{-{}-save PATH}, \code{-{}-csv PATH} & save the algorithm
      (\code{Algorithm.save}) and the training episodes
      (\code{TrainingLog.to\_csv}) \\
    \quad \code{-{}-threads N} & PyTorch threads [1]; 0 keeps the PyTorch
      default \\
    \code{compare ENV\_ID} & train baseline algorithms over one or more
      seeds and compare them with random actions and the classical controller
      on the same episodes (Section~\ref{sec:marl-baselines}) \\
    \quad \code{-{}-algos ALGO ...} & algorithms [every one with a preset for
      the environment] \\
    \quad \code{-{}-seeds S ...} & training seeds [0] \\
    \quad \code{-{}-steps N}, \code{-{}-num-envs N} & budget and copies per run
      [preset] \\
    \quad \code{-{}-episodes N}, \code{-{}-eval-seed S} & evaluation episodes
      [64] and seed [1] \\
    \quad \code{-{}-set}, \code{-{}-env-kwarg} & as for \code{run} \\
    \quad \code{-{}-markdown}, \code{-{}-csv PATH} & print a Markdown table,
      save the report as CSV \\
    \quad \code{-{}-save-dir DIR} & save every trained algorithm as
      \code{DIR/NAME\_seedS.pt} \\
    \code{evaluate CHECKPOINT ENV\_ID} & evaluate a saved algorithm against
      random actions and the baseline \\
    \quad \code{-{}-episodes N}, \code{-{}-seed S} & evaluation episodes [64]
      and seed [1] \\
    \quad \code{-{}-env-kwarg KEY=VALUE ...} & environment arguments (as used
      in training) \\
    \bottomrule
  \end{tabularx}
\end{table}

\begin{shellcode}
marl-train list                       # algorithms and references
marl-train presets                    # tuned (algorithm, environment) pairs
marl-train run mappo PowerGrid-v0     # train with the preset, then evaluate
marl-train run qmix LineMsg-v0 --steps 50000 --set lr=1e-3
marl-train run maddpg Consensus-v0 --save runs/maddpg.pt --csv runs/maddpg.csv
marl-train evaluate runs/maddpg.pt Consensus-v0 --episodes 64
marl-train run ippo Formation-v0 --no-preset --steps 200000 --num-envs 32 \
    --set lr=1e-3 normalize_observations=False --env-kwarg formation_shape=wedge
\end{shellcode}

The evaluation of \code{run} and \code{evaluate} uses
\code{make\_vector\_env(env\_id, min(episodes, 64))}, so up to 64 episodes run in
parallel. Output of \code{marl-train run qmix LineMsg-v0} (the complete preset;
wall-clock times depend on the machine):

\begin{textcode}
Training qmix on LineMsg-v0 (preset): 25,000 env steps, 16 copies, seed 0
       2,512 steps  mean return (last 20)        43.02     2.7 s
       5,008 steps  mean return (last 20)        54.33     4.7 s
       ...
      22,512 steps  mean return (last 20)        92.35    17.3 s
      25,008 steps  mean return (last 20)        91.96    19.1 s
qmix on LineMsg-v0: 496 episodes, 25008 env steps, mean return (last 100) 92.18, 17.0 s

Evaluation on LineMsg-v0, 64 episodes (seed 1):
  policy       mean return   95% CI +-   length
  random              43.2        1.17     50.0
  qmix                  95           0     50.0
  baseline              95           0     50.0
\end{textcode}

The progress lines report the mean return of the last 20 exploratory training
episodes; the summary line is \code{TrainingLog.summary()}, whose step count
and time refer to the last record of the log (an update or a finished
episode). The evaluation columns
are the mean team return of the deterministic policy, the half-width of its
95\,\% confidence interval and the mean episode length
(Table~\ref{tab:workflow-result}).

% ---------------------------------------------------------------------------
\section{Using the algorithms as baselines}\label{sec:marl-baselines}

A new method for one of the environments is judged against established ones.
Three kinds of reference are available for every environment: uniformly
random actions (the floor), the environment's classical controller
(\code{env\_lib.baseline\_policy}, the standard of its field) and the seven
learning algorithms, with tuned presets for the environments of
Table~\ref{tab:marl-presets}. The function \code{marl\_algorithms.compare}
trains the learning baselines, over as many seeds as you ask for, and
evaluates everything, your methods included, on the same seeded episodes.

\subsection{In one call}

\begin{pycode}
import numpy as np
from marl_algorithms import compare

def my_policy(obs):
    """Your method: (num_envs, n_agents, obs_dim) -> (num_envs, n_agents)."""
    return np.zeros(obs.shape[:-1], dtype=np.float32)   # replace me

report = compare("PowerGrid-v0", ["mappo", "ippo"], seeds=(0, 1),
                 policies={"my method": my_policy},
                 total_steps=4_096, num_envs=16)       # small budget: seconds
print(report)
report.to_csv("runs/power_grid_comparison.csv")
\end{pycode}

Without \code{total\_steps} and \code{num\_envs} every algorithm trains with
its preset. The script \code{examples/algorithms/baseline\_comparison.py} runs
exactly this comparison with a proposed method, a distributed droop controller
that also reacts to the neighbours' mean frequency deviation, and the full
presets (about six minutes on one core):

\begin{textcode}
PowerGrid-v0: mean team return over 64 evaluation episodes (seed 1); higher is better

method     kind       mean return    std  seeds  env steps  train [s]
---------  ---------  -----------  -----  -----  ---------  ---------
random     reference       -629.3      -      -          -          -
baseline   reference       -0.786      -      -          -          -
mappo      algorithm        -1.02  0.182      2    251,904      101.9
ippo       algorithm        -2.25  0.172      2    200,704       65.2
my method  policy          -0.675      -      -          -          -
\end{textcode}


\begin{figure}[htbp]
  \centering
  \includegraphics[width=0.72\textwidth]{baseline_comparison.png}
  \caption{Output figure of \code{baseline\_comparison.py}: mean return on
    \envid{PowerGrid-v0} over 64 evaluation episodes; the error bars are the
    standard deviation over two training seeds. Random actions ($-629.3$) are
    off the scale.}
  \label{fig:marl-baselines}
\end{figure}

Figure~\ref{fig:marl-baselines} draws the same numbers. The rows are the
references first, then the algorithms with the mean and the
standard deviation over the training seeds, the environment steps and the
wall-clock training time of one run, then your policies. \code{compare} works
in four steps:

\begin{enumerate}
  \item it creates the evaluation environment,
    \code{make\_vector\_env(env\_id, min(n\_episodes, 64), **env\_kwargs)};
  \item it evaluates random actions, the classical controller (when the
    environment has one) and your policies, first, so that an error in a
    policy shows before any training;
  \item it trains every algorithm once per seed, with its preset
    (\code{train\_preset}) or, without a preset, with its default
    configuration and \code{total\_steps}, and evaluates the deterministic
    policy;
  \item it returns a \code{Comparison}, which also keeps every trained
    algorithm and training log.
\end{enumerate}

Every evaluation resets the same environment copies with the same seed, so all
rows see identical initial conditions and disturbances. With the default
training seed 0 and evaluation seed 1 the numbers reproduce the published
results of Table~\ref{tab:marl-results-table}.

\subsection{From the command line}

\begin{shellcode}
marl-train compare PowerGrid-v0                       # every preset algorithm, seed 0
marl-train compare PowerGrid-v0 --algos mappo ippo --seeds 0 1 2 --csv runs/pg.csv
marl-train compare LineMsg-v0 --markdown              # a table for a README or a paper
marl-train compare Consensus-v0 --save-dir runs/baselines  # keep the trained baselines
marl-train compare Formation-v0 --algos mappo --steps 200000   # no preset: set a budget
\end{shellcode}

\code{marl-train compare} checks the environment, the algorithm names, their
action kinds and the \code{-{}-set} fields before training
(Table~\ref{tab:marl-cli}).

\subsection{Recipes}

\begin{description}[leftmargin=0.6cm,style=nextline]
  \item[Choose the baselines and the budget]
    \code{algorithms=["mappo", "maddpg"]} selects methods (by default every
    algorithm with a preset for the environment, \code{marl-train presets});
    \code{seeds=(0, 1, 2, 3, 4)} adds seeds; \code{total\_steps} and
    \code{num\_envs} replace the presets' budget, for example for a quick
    check before a long run.
  \item[Another configuration of the environment]
    \code{env\_kwargs=\{"n\_buses": 32\}} trains and evaluates every method on
    the variant. The presets' hyperparameters are kept; a larger variant may
    need a larger \code{total\_steps}.
  \item[An algorithm without a preset]
    Pass \code{total\_steps} to train it with its default configuration, or
    tune it yourself with \code{train(...)} and pass the trained algorithm in
    \code{policies}.
  \item[A policy written for a single environment]
    \code{compare} evaluates batched observations. Wrap a policy that maps one
    observation \code{(n\_agents, obs\_dim)} to one action with
    \code{per\_copy(policy)}; batched policies are faster.
  \item[A method built on this package]
    A modified algorithm (a new loss, another critic) is an \code{Algorithm}:
    train it and pass the object in \code{policies}; its deterministic policy
    is evaluated like the baselines. For several seeds of your method, train
    one per seed and pass them under separate names, or register it
    (Section~\ref{sec:marl-extending}) and list it in \code{algorithms}.
  \item[Reuse the trained baselines]
    \code{report.algorithms[("mappo", 0)]} is the trained MAPPO of seed 0:
    save it, record an episode with \code{record\_episode(env,
    algo.policy(), path)} or evaluate it on other scenarios;
    \code{report.logs} holds the training logs for learning curves
    (Section~\ref{sec:plotkit-curves}).
  \item[Tables and further analysis]
    \code{report.to\_markdown()}, \code{report.to\_csv(path)} and
    \code{report.records()} (one dictionary per row, for example for
    \code{pandas.DataFrame}); \code{report["mappo"].returns} gives the
    per-seed returns and \code{report.ranking()} orders the algorithms and
    your policies by mean return.
\end{description}

% ---------------------------------------------------------------------------
\section{Design of the core}\label{sec:marl-core}

The package separates what all methods share, what distinguishes each
method, and the settings tuned for the environments. Every algorithm has its
own module, next to the \code{common} module of its family, and the algorithm
modules contain only what is specific to them:

\begin{textcode}
src/marl_algorithms/
  __init__.py           lazy exports (train, compare, algorithms, ...)
  __main__.py           the marl-train command line
  registry.py           AlgorithmInfo, get_algorithm, make_algorithm, train
  baselines.py          compare, Comparison, per_copy: baselines
  core/                 shared by all methods
    spec.py             MultiAgentSpec: per-agent view of the joint spaces
    runner.py           make_vector_env, VectorRunner, Transition
    buffers.py          RolloutBuffer, compute_gae, ReplayBuffer
    normalization.py    RunningMeanStd, ObservationNormalizer, RewardScaler
    networks.py         mlp, policies, QMixer, PerAgent, agent_one_hot,
                        soft_update
    config.py           AlgorithmConfig, OnPolicyConfig, OffPolicyConfig
    base.py             Algorithm, OnPolicyAlgorithm, OffPolicyAlgorithm,
                        TrainingLog
  algorithms/           one module per algorithm, grouped by family
    ppo/                on-policy
      common.py         PPOConfig and the shared PPO implementation
      ippo.py           IPPO (decentralised critics)
      mappo.py          MAPPO (centralised critic)
    ddpg/               off-policy actor-critic
      common.py         DDPGConfig and the shared implementation
      maddpg.py         MADDPG
      matd3.py          MATD3 (twin critics, smoothing, delayed actors)
    q_learning/         value-based
      common.py         QLearningConfig and the shared implementation
      iql.py            IQL (no mixing)
      vdn.py            VDN and VDNMixer (additive mixing)
      qmix.py           QMIX (monotonic mixing with QMixer)
  presets/              tuned settings for the env_lib environments
    __init__.py         PRESETS, get_preset, list_presets, train_preset
    ppo.py, ddpg.py, q_learning.py
                        the tables of each family, with tuning notes
\end{textcode}

The docstring of every family package (for example
\code{help(marl\_algorithms.algorithms.ppo)}) describes the methods, the
implementation choices and the references; the docstring of every class gives
its losses.

\subsection{The per-agent view: \texorpdfstring{\code{MultiAgentSpec}}{MultiAgentSpec}}\label{sec:marl-spec}

\code{MultiAgentSpec.from\_env(env)} reads the agent structure of a single or
vector environment (from \code{single\_observation\_space} and
\code{single\_action\_space} when present) and
\code{MultiAgentSpec.from\_spaces(observation\_space, action\_space)} from the
spaces. The spec is a frozen dataclass with the fields \code{n\_agents},
\code{obs\_dim}, \code{action\_kind} (\code{"continuous"} or
\code{"discrete"}), \code{action\_dim} (1 for discrete actions),
\code{n\_actions} (discrete only), \code{action\_low} and \code{action\_high}
(continuous only, shape \code{(n\_agents, action\_dim)}), \code{obs\_shape},
\code{action\_shape} and \code{action\_dtype}. The observation space must be a
\code{Box}; the action spaces are read as follows.

\begin{itemize}
  \item \code{Box} of shape \code{(n, d)} with observations \code{(n, obs\_dim)}:
    $n$ agents with $d$-dimensional continuous actions; shape \code{(n,)}: $n$
    agents with one action each. Any other \code{Box} is a single agent whose
    action is the flattened box (the Kuramoto environments). The bounds must be
    finite; they may differ between agents and dimensions.
  \item \code{MultiDiscrete} and \code{MultiBinary}: one agent per entry, with
    the same number of choices for every agent; \code{MultiDiscrete} must
    start at 0.
  \item \code{Discrete}: only for a single agent. A joint \code{Discrete}
    action over several agents, such as the default \code{Discrete(2 ** n)}
    action of LineMsg, raises \code{TypeError}; \code{make\_vector\_env}
    therefore creates LineMsg with \code{action\_space\_type="multibinary"}.
  \item With several agents the observation must be stacked as
    \code{(n, obs\_dim)}; with one agent it is flattened.
\end{itemize}

\code{spec.agent\_obs(obs)} returns observations with any leading batch shape
as \code{float32} arrays \code{(*batch, n\_agents, obs\_dim)};
\code{spec.env\_action(actions)} converts per-agent actions,
\code{(*batch, n\_agents, action\_dim)} floats (or \code{(*batch, n\_agents)}
when \code{action\_dim} is 1) or \code{(*batch, n\_agents)} integers, into the
joint action and clips continuous actions to the bounds. The property
\code{state\_dim} $= n \cdot \code{obs\_dim}$ is the size of the global state
$s = [o_1, \dots, o_n]$ that the centralised critics and the QMIX mixer use:
none of the environments exposes a separate state, so the global state is the
concatenation of all agents' observations. \code{describe()} returns a one-line
summary such as \code{"16 agents, 10-d observations, continuous 1-d per agent"}.

\subsection{Collecting experience: \texorpdfstring{\code{VectorRunner}}{VectorRunner} and \texorpdfstring{\code{Transition}}{Transition}}\label{sec:marl-runner}

% norun
\begin{pycode}
make_vector_env(env_id, num_envs=16, **env_kwargs)
# = env_lib.make_vec(env_id, num_envs, autoreset_mode="same_step", **env_kwargs)
\end{pycode}
creates the training copies; for LineMsg it sets
\code{action\_space\_type="multibinary"} unless the argument is given. Environments without a native vector implementation (LineMsg,
WirelessComm, Pistonball, AJLATT and the PyTorch Kuramoto ids) are simulated
with Gymnasium's \code{SyncVectorEnv}.

\code{VectorRunner(envs, spec=None)} steps such a vector environment in
per-agent form; it raises \code{ValueError} for any autoreset mode other than
\code{"same\_step"}. \code{reset(seed=None)} resets all copies and returns
observations \code{(B, n, obs\_dim)}; \code{step(actions)} applies per-agent
actions to all copies and returns a \code{Transition}
(Table~\ref{tab:marl-transition}), after which \code{runner.obs} holds the
observations to act on next; \code{pop\_episodes()} returns and clears the
\code{(return, length)} pairs of the episodes finished so far.

\begin{table}[htbp]
  \centering\small
  \caption{Fields of \code{Transition} ($B$ copies, $n$ agents).}
  \label{tab:marl-transition}
  \begin{tabularx}{\textwidth}{@{}>{\raggedright\arraybackslash}p{3.3cm}>{\raggedright\arraybackslash}p{3.4cm}>{\raggedright\arraybackslash}X@{}}
    \toprule
    Field & Shape & Content \\
    \midrule
    \code{obs} & \code{(B, n, obs\_dim)} & observations the actions were taken in \\
    \code{actions} & \code{(B, n, action\_dim)} or \code{(B, n)} & per-agent actions as passed to \code{step} (unclipped) \\
    \code{reward} & \code{(B,)} & team rewards \\
    \code{agent\_rewards} & \code{(B, n)} & per-agent rewards from \code{info["agent\_rewards"]}; the team reward for every agent when the environment reports none \\
    \code{next\_obs} & \code{(B, n, obs\_dim)} & successor observations; the \emph{final} observation for copies whose episode ended \\
    \code{terminated}, \code{truncated} & \code{(B,)} & episode-end flags; the property \code{done} is their disjunction \\
    \bottomrule
  \end{tabularx}
\end{table}

With same-step autoreset a finished copy is reset within the step, so the
observation returned by the environment already belongs to the next episode.
The runner therefore takes the final observation from
\code{infos["final\_obs"]} and the final per-agent rewards from
\code{infos["final\_info"]["agent\_rewards"]}. Every transition is a real one
(there is no reset step to skip, as with next-step autoreset), and the value
of the final observation is available to bootstrap episodes that were
truncated by the time limit. Only \code{terminated} stops the bootstrap in the
algorithms; \code{truncated} only cuts the advantage recursion.

\subsection{Buffers and generalised advantage estimation}\label{sec:marl-buffers}

\paragraph{Rollouts.} \code{RolloutBuffer(length, num\_envs, fields)} stores
named arrays of shape \code{(T, B, ...)}, with \code{fields} a dictionary
\code{\{name: (trailing\_shape, dtype)\}}. \code{add(**values)} stores one
vector step (every field must be given), \code{buffer[name]} returns the
stored steps, \code{full} tells whether $T$ steps are stored and
\code{reset()} starts a new rollout in the same arrays.

\paragraph{GAE.} The function

% norun
\begin{pycode}
compute_gae(rewards, values, next_values, terminated, truncated, gamma,
            gae_lambda) -> (advantages, returns)
\end{pycode}
works on arrays of shape \code{(T, B, ...)}; trailing dimensions such as
agents are broadcast against the \code{(T, B)} flags. With $V_t = V(o_t)$,
$V'_t$ the value of the successor observation $o'_t$ (the final observation
when the episode ended), $\tau_t$ the \code{terminated} flag of step $t$ and
$\mathrm{done}_t$ the flag that is set when the episode terminated or was
truncated at step $t$,
\begin{align}
  \delta_t &= r_t + \gamma\,(1 - \tau_t)\, V'_t - V_t, \label{eq:marl-td}\\
  A_t &= \delta_t + \gamma \lambda\, (1 - \mathrm{done}_t)\, A_{t+1},
  \qquad A_T = 0, \label{eq:marl-gae}\\
  R_t &= A_t + V_t. \nonumber
\end{align}
The recursion runs backwards from the last step of the rollout, where the
bootstrap enters through $V'_{T-1}$ in $\delta_{T-1}$. A terminated episode
does not bootstrap, a truncated one bootstraps from the value of its final
observation, and neither passes advantages across the episode boundary. A test
compares the function with this recursion written out step by step.

\paragraph{Replay.} \code{ReplayBuffer(capacity, spec)} is a circular buffer of
joint transitions. Its method \code{add(transition)} stores every copy of a
vector transition, one entry per copy and step, so the capacity counts
transitions rather than vector steps. The method \code{sample(batch\_size,
rng)} draws transitions uniformly with replacement using the algorithm's
generator and returns a dictionary of NumPy arrays:

\begin{textcode}
obs, next_obs      (N, n, obs_dim)       float32
actions            (N, n, action_dim)    float32   (continuous actions)
                   (N, n)                int64     (discrete actions)
reward             (N,)                  float32
agent_rewards      (N, n)                float32
terminated         (N,)                  float32
\end{textcode}

Truncation is not stored, because the stored successor of a truncated
transition is its final observation and the target bootstraps from it.

\subsection{Normalisation}\label{sec:marl-normalisation}

\code{RunningMeanStd(shape=())} keeps a running mean $\mu$, variance
$\sigma^2$ and count $N$ (initially $10^{-4}$) and includes a batch with mean
$\mu_b$, variance $\sigma_b^2$ and size $N_b$ with the parallel update rule
\begin{equation}
  \Delta = \mu_b - \mu, \quad
  \mu \leftarrow \mu + \Delta \frac{N_b}{N + N_b}, \quad
  \sigma^2 \leftarrow \frac{N \sigma^2 + N_b \sigma_b^2
      + \Delta^2 N N_b / (N + N_b)}{N + N_b}, \quad
  N \leftarrow N + N_b .
\end{equation}

\code{ObservationNormalizer(obs\_dim, clip=10.0, enabled=True)} standardises
observations feature by feature, $\hat o = \operatorname{clip}((o - \mu) /
\sigma, -10, 10)$, with statistics \emph{shared by all agents}: every agent's
observation of every copy is a sample of the same per-feature statistics, as
befits homogeneous agents with shared parameters. \code{update(obs)} adds
observations of any leading shape; with \code{enabled=False} the normaliser
returns its input unchanged.

\code{RewardScaler(num\_envs, gamma, enabled=True)} divides rewards by the
running standard deviation of the discounted return. Every copy $b$ keeps a
return $G_b \leftarrow \gamma G_b + \bar r_b$, where $\bar r_b$ is the reward of
copy $b$ averaged over its trailing dimensions (the agents), the statistics
include the new $G_b$, the returns of finished copies are reset to zero, and
the rewards are returned as $r / \sigma_G$. Rewards are scaled, not centred, so
the optimal policy is unchanged, while value targets become of order one
whether an environment's rewards are of order $10^{-3}$ per step (PowerGrid)
or $10^{2}$ (Consensus).

\subsection{Networks}\label{sec:marl-networks}

The building blocks of \pkg{marl\_algorithms.core.networks} take per-agent
inputs with arbitrary leading dimensions \code{(..., features)}:

% norun
\begin{pycode}
mlp(in_dim, out_dim, hidden_sizes=(64, 64), activation="tanh", *,
    output_gain=1.0, layer_norm=False)
PerAgent(factory, n_agents, shared=True)
GaussianPolicy(in_dim, action_dim, low, high, hidden_sizes, activation,
               log_std_init=-0.5)
CategoricalPolicy(in_dim, n_actions, hidden_sizes, activation)
DeterministicPolicy(in_dim, action_dim, low, high, hidden_sizes,
                    activation="relu")
QMixer(n_agents, state_dim, embed_dim=32, hypernet_hidden=64)
agent_one_hot(batch_shape, n_agents, device)
soft_update(target, source, tau)
\end{pycode}

\begin{itemize}
  \item \code{mlp} builds a multi-layer perceptron with orthogonal
    initialisation (gain $\sqrt{2}$ for ReLU-type activations, $5/3$ for
    tanh) and zero biases; the activation is \code{"tanh"}, \code{"relu"},
    \code{"elu"} or \code{"gelu"}. The policy heads use
    \code{output\_gain=0.01}, which makes the initial policies nearly uniform
    (categorical) or centred in the action range (Gaussian, deterministic).
  \item \code{PerAgent} with \code{shared=True} applies one network, built by
    \code{factory()}, to all agents in one call; with \code{shared=False}
    every agent has its own network and agent $i$'s slice of the input goes
    through network $i$. \code{call(method, x, *args)} dispatches any method
    of the underlying networks in the same way (tensor arguments are sliced
    per agent, others, such as a \code{torch.Generator}, are passed
    unchanged).
  \item Agent identifiers: when parameters are shared and
    \code{agent\_ids=True} (both defaults) and there is more than one agent,
    the one-hot identifier $e_i \in \{0, 1\}^n$ of agent $i$ is appended to
    its input, $x_i = [\hat o_i, e_i]$, so that agents sharing one network can
    still behave differently. \code{Algorithm.agent\_inputs(obs)} builds these
    inputs for all agents at once with \code{agent\_one\_hot}.
  \item \code{GaussianPolicy} is a diagonal Gaussian with a state-independent
    standard deviation, expressed in units of the action range: with the
    centre $c = (\text{high} + \text{low})/2$ and the half-range
    $h = (\text{high} - \text{low})/2$,
    $\pi_\theta(\cdot \mid x) = \mathcal{N}\bigl(c + h \odot f_\theta(x),\;
    \operatorname{diag}(h \odot e^{\ell})^2\bigr)$ with the learned vector
    $\ell$ initialised to \code{log\_std\_init}. Samples are unbounded; the
    log-density $\sum_k \log \mathcal{N}(a_k)$ is that of the unclipped sample,
    and the entropy is $\sum_k \bigl(\tfrac12 + \tfrac12 \log 2\pi + \log
    \sigma_k\bigr)$. A shared policy keeps the bounds of all agents, so agents
    with different action ranges can share it.
  \item \code{CategoricalPolicy} returns the logits of a categorical
    distribution; like \code{GaussianPolicy} it provides
    \code{sample(x, generator)} (actions and log-probabilities) and
    \code{log\_prob(x, action)} (log-probabilities and entropy).
  \item \code{DeterministicPolicy} computes $\mu_\theta(x) = c + h \odot
    \tanh f_\theta(x)$, always within the bounds.
  \item \code{QMixer} is the monotonic mixing network of QMIX
    (Section~\ref{sec:marl-value}), and \code{soft\_update} performs Polyak
    averaging, $\theta' \leftarrow (1 - \tau)\theta' + \tau\theta$.
\end{itemize}

\subsection{Base classes and training loops}\label{sec:marl-base}

Every algorithm derives from \code{Algorithm}:

% norun
\begin{pycode}
Algorithm(spec, config=None, *, device="cpu", seed=None, **overrides)
\end{pycode}

\code{spec} is a \code{MultiAgentSpec} or an environment (single or vector) to
read it from; \code{config} an instance of the class attribute
\code{config\_class} (defaults to \code{config\_class()}); \code{overrides}
replace configuration fields with \code{dataclasses.replace}. An environment
whose action kind is not in the class attribute \code{action\_kinds}, a
configuration of the wrong class and an unknown field raise \code{TypeError};
invalid values raise \code{ValueError} when the configuration is created.
Table~\ref{tab:marl-config-common} lists the fields shared by all algorithms.

\begin{table}[htbp]
  \centering\small
  \caption{Configuration fields shared by the algorithms
    (\pkg{marl\_algorithms.core.config}). \code{PPOConfig} derives from
    \code{OnPolicyConfig}, \code{DDPGConfig} and \code{QLearningConfig} from
    \code{OffPolicyConfig}; both derive from \code{AlgorithmConfig}.}
  \label{tab:marl-config-common}
  \begin{tabularx}{\textwidth}{@{}>{\raggedright\arraybackslash}p{4.3cm}>{\raggedright\arraybackslash}p{1.8cm}>{\raggedright\arraybackslash}X@{}}
    \toprule
    Field & Default & Meaning \\
    \midrule
    \multicolumn{3}{@{}l}{\emph{\code{AlgorithmConfig} (all algorithms)}} \\
    \code{hidden\_sizes} & \code{(64, 64)} & widths of the hidden layers of every network \\
    \code{activation} & \code{"tanh"} & hidden activation (\code{"tanh"}, \code{"relu"}, \code{"elu"}, \code{"gelu"}) \\
    \code{gamma} & 0.99 & discount factor $\gamma$ \\
    \code{share\_parameters} & \code{True} & one set of network weights for all agents \\
    \code{agent\_ids} & \code{True} & append one-hot agent identifiers to per-agent inputs when parameters are shared \\
    \code{max\_grad\_norm} & 10.0 & gradient-norm clipping threshold \\
    \midrule
    \multicolumn{3}{@{}l}{\emph{\code{OnPolicyConfig} (IPPO, MAPPO)}} \\
    \code{max\_grad\_norm} & 0.5 & overrides the default above \\
    \code{rollout\_length} & 64 & vector steps $T$ per rollout ($T \cdot B$ transitions per update) \\
    \code{lr} & $3 \cdot 10^{-4}$ & Adam learning rate (of the actor for PPO) \\
    \code{gae\_lambda} & 0.95 & GAE parameter $\lambda$ \\
    \code{normalize\_observations} & \code{True} & standardise observations with running statistics \\
    \code{normalize\_rewards} & \code{True} & scale rewards by the running standard deviation of the return \\
    \midrule
    \multicolumn{3}{@{}l}{\emph{\code{OffPolicyConfig} (MADDPG, MATD3, IQL, VDN, QMIX)}} \\
    \code{activation} & \code{"relu"} & overrides the default above \\
    \code{buffer\_size} & 200\,000 & replay capacity in transitions (one per copy and step) \\
    \code{batch\_size} & 256 & transitions per gradient step \\
    \code{warmup\_steps} & 2\,000 & environment steps (summed over copies) with uniformly random actions before learning starts \\
    \code{update\_every} & 1 & vector steps between rounds of gradient steps \\
    \code{gradient\_steps} & 1 & gradient steps per round \\
    \code{tau} & 0.01 & Polyak coefficient $\tau$ of the target networks \\
    \code{reward\_scale} & 1.0 & constant factor $\kappa$ applied to rewards before learning \\
    \bottomrule
  \end{tabularx}
\end{table}

\paragraph{Seeding without global side effects.} The seed initialises a
\code{numpy.random.SeedSequence}, from which the algorithm derives its own
NumPy generator \code{self.np\_rng} (random warm-up actions, replay sampling,
$\epsilon$-greedy draws) and a seed for its own \code{torch.Generator}
\code{self.torch\_rng} (policy samples, exploration and target-smoothing noise,
minibatch permutations). The networks are built inside
\code{torch.random.fork\_rng}, seeded from the same sequence, so network
initialisation is reproducible while PyTorch's global random state is restored
afterwards. Apart from this temporary, forked seeding, no code of the package
seeds or draws from the global NumPy or PyTorch generators, and none changes
the number of PyTorch threads; the scripts and \code{marl-train} set the
thread count themselves. With the same seed and the same machine, training is identical
(the tests check this for every algorithm); \code{seed=None} draws fresh
entropy.

\paragraph{The on-policy loop.} \code{OnPolicyAlgorithm.learn} resets the
copies with the given seed and then repeats until \code{env\_steps} reaches
\code{total\_steps}:
\begin{enumerate}
  \item for \code{rollout\_length} vector steps, call
    \code{\_rollout\_step(obs)} for exploratory actions and extra fields (for
    PPO the normalised observations and the log-probabilities), step the
    runner, pass the transition to \code{\_observe\_transition} and store it
    with the extra fields in a \code{RolloutBuffer};
  \item call \code{\_update(buffer)} and record its statistics and the
    finished episodes in the log, then call the callback.
\end{enumerate}
The buffer always holds the fields \code{obs}, \code{actions}, \code{reward},
\code{agent\_rewards}, \code{next\_obs}, \code{terminated} and
\code{truncated}, plus the fields declared by \code{\_rollout\_fields()}.

\paragraph{The off-policy loop.} \code{OffPolicyAlgorithm.learn} resets the
copies and then, at every vector step, takes uniformly random actions within
the bounds while \code{env\_steps < warmup\_steps} and \code{\_explore(obs)}
afterwards, steps the runner and adds the transition to the replay buffer
(created on the first call), and records the episodes finished in this step
in the log. Once the warm-up is over and the buffer holds at least
\code{batch\_size} transitions, every \code{update\_every} vector steps it
performs \code{gradient\_steps} calls of \code{\_update(batch)} on batches
sampled from the replay buffer, records the statistics of the last one and
calls the callback. With $B$ copies the loop therefore performs
\code{gradient\_steps} / (\code{update\_every} $\cdot B$) gradient steps per
environment step; the presets use one gradient step per vector step of 16
copies.

\paragraph{Public methods.} Every algorithm provides the following methods
and the counters \code{env\_steps} and \code{num\_updates}:

% norun
\begin{pycode}
algo.learn(envs, total_steps, *, seed=None, log=None,
           callback=None) -> TrainingLog
algo.act(obs, deterministic=False)   # (*batch, n, obs_dim) -> agent actions
algo.policy(deterministic=True)      # env observation -> env action
algo.evaluate(env, n_episodes=10, *, seed=None, deterministic=True)
algo.state_dict(); algo.load_state_dict(state)
algo.save(path) -> Path
Algorithm.load(path, device="cpu") -> Algorithm      # class method
\end{pycode}
The method \code{evaluate} runs \code{env\_lib.evaluate} with the policy
returned by \code{policy} and returns its \code{EvaluationResult}.

\paragraph{The training log.} \code{TrainingLog} records every finished
training episode as \code{(env\_steps, team\_return, length)} in
\code{episodes} and the statistics of every update as
\code{(env\_steps, statistics)} in \code{updates}; the returns are those of
the environment's team reward, before any reward scaling.
\code{episode\_steps} and \code{episode\_returns} return them as arrays,
\code{mean\_return(last=100)} averages the last episodes (NaN before the
first), \code{curve(points=50)} bins the returns into equal intervals of
environment steps for a learning curve, \code{summary()} gives a one-line
summary, \code{to\_csv(path)} writes the episodes as CSV
(\code{env\_steps, return, length}), and \code{wall\_time} is the time from the
creation of the log to its last record. Section~\ref{sec:plotkit-curves} shows how to plot learning
curves of several seeds with confidence bands.

% ---------------------------------------------------------------------------
\section{PPO: IPPO and MAPPO}\label{sec:marl-ppo}

IPPO and MAPPO train decentralised stochastic actors $\pi_\theta(a_i \mid x_i)$
with the clipped surrogate objective of PPO (Schulman et al., 2017) and
generalised advantage estimation (Schulman et al., 2016), and differ only in
the critic:

\begin{itemize}
  \item \textbf{IPPO} gives every agent a decentralised critic
    $V_\phi(x_i)$ on its own input $x_i = [\hat o_i, e_i]$, trained on the
    agent's own reward $r^i_t$ (\code{reward\_source="agent"}, the default of
    IPPO). Every agent is an independent PPO learner that treats the others as
    part of the environment.
  \item \textbf{MAPPO} gives every agent a centralised critic
    $V_\phi(s, i)$ on the global state $s = [\hat o_1, \dots, \hat o_n]$, the
    concatenation of all normalised observations, trained on the team reward
    $r_t$ (\code{reward\_source="team"}, the default of MAPPO). With shared
    parameters one network takes $[s, e_i]$, which amounts to one value head
    per agent over the shared state; without sharing agent $i$ has its own
    network on $s$. The critic is used only in training, so execution stays
    decentralised.
\end{itemize}

\code{reward\_source} selects the reward of the critics for both methods:
\code{"team"} gives every critic the team reward, \code{"agent"} gives critic
$i$ agent $i$'s own reward. The presets of MAPPO on \envid{Platoon-v0} and
\envid{Consensus-v0} use \code{"agent"} (Section~\ref{sec:marl-results}).

\paragraph{Acting.} In a rollout the observation normaliser first includes the
new observations and then normalises them; the normalised observations are
stored, so the update evaluates the policy on exactly the input that produced
each action. Continuous actions are sampled unbounded and stored unclipped
(the environment adapter clips them); the stored log-probability is that of
the unclipped sample. The deterministic policy returns the mean, clipped to
the bounds, or the most likely discrete action.

\paragraph{Advantages and value targets.} After a rollout of $T$ vector steps
on $B$ copies the critic values $V_t^i$ of the stored inputs and $V'^i_t$ of
the successor observations (normalised with the current statistics) give, for
every agent and with the rewards $r^i_t$ of the chosen source scaled by the
\code{RewardScaler}, the GAE advantages $A^i_t$ and targets $R^i_t = A^i_t +
V^i_t$ of \eqref{eq:marl-td} and \eqref{eq:marl-gae}. With
\code{normalize\_advantages} the advantages are standardised over the whole
rollout (all steps, copies and agents), $\hat A = (A - \bar A)/(\operatorname{std}
A + 10^{-8})$.

\paragraph{Losses.} With the probability ratio
$\rho^i_t = \pi_\theta(a^i_t \mid x^i_t) / \pi_{\theta_\text{old}}(a^i_t \mid
x^i_t)$ and the means taken over the samples of a minibatch and all agents,
the actor minimises
\begin{equation}
  L^{\pi}(\theta) = -\operatorname{mean}\Bigl[\min\bigl(\rho^i_t \hat A^i_t,\;
    \operatorname{clip}(\rho^i_t, 1 - \epsilon, 1 + \epsilon)\, \hat A^i_t\bigr)\Bigr]
    - c_H \operatorname{mean}\bigl[\mathcal{H}\bigl(\pi_\theta(\cdot \mid x^i_t)\bigr)\bigr]
  \label{eq:marl-ppo-actor}
\end{equation}
with $\epsilon$ = \code{clip\_range} and $c_H$ = \code{ent\_coef}, and the
critic minimises
\begin{equation}
  L^{V}(\phi) = \operatorname{mean}\Bigl[\max\Bigl(\ell\bigl(V_\phi - R\bigr),\;
    \ell\bigl(V_\text{old} + \operatorname{clip}(V_\phi - V_\text{old},
    -\epsilon_V, \epsilon_V) - R\bigr)\Bigr)\Bigr],
  \label{eq:marl-ppo-critic}
\end{equation}
where $V_\phi$ is $V_\phi(x^i_t)$ for IPPO and $V_\phi(s_t, i)$ for MAPPO,
$V_\text{old}$ the value before the update, $\epsilon_V$ =
\code{value\_clip\_range} (\code{None} removes the clipped term) and
$\ell(e) = e^2/2$, or with \code{huber\_delta} $=\eta$ the Huber loss
$\ell(e) = e^2/2$ for $|e| \le \eta$ and $\eta(|e| - \eta/2)$ otherwise.

\paragraph{Optimisation.} A sample is one copy at one step with all its agents,
because the centralised critic needs every agent's observation; a rollout
therefore holds $T B$ samples. Each of the \code{n\_epochs} epochs shuffles
them with the algorithm's generator and splits them into
\code{num\_minibatches} minibatches; on each minibatch the actor takes one Adam
step on \eqref{eq:marl-ppo-actor} and then the critic one Adam step on
\eqref{eq:marl-ppo-critic}, with separate optimisers ($\epsilon_\text{Adam} =
10^{-5}$, learning rates \code{lr} and \code{critic\_lr}) and separate
gradient-norm clipping at \code{max\_grad\_norm}. With \code{target\_kl} the
epochs stop early once the mean approximate KL divergence of an epoch,
$\operatorname{mean}[(\rho - 1) - \log \rho]$, exceeds $1.5$ times
\code{target\_kl}. With \code{anneal\_lr} both learning rates decay linearly to
zero over the steps of the current \code{learn} call. Every update records
\code{policy\_loss}, \code{value\_loss}, \code{entropy}, \code{approx\_kl},
\code{clip\_fraction} (the fraction of $|\rho - 1| > \epsilon$), averaged over
the minibatches, and \code{explained\_variance}, \code{epochs},
\code{learning\_rate}, \code{value\_mean}, \code{action\_std} (continuous
actions) and \code{reward\_std} (with reward scaling).

\paragraph{Stated simplifications.} The implementation follows the PPO recipe
of the references (including the practices studied by Andrychowicz et al.,
2021) with four standard simplifications, stated in the module docstring:
feed-forward instead of recurrent networks (the \pkg{env\_lib} observations
carry the recent history that matters, such as previous actions and rates of
change); a state-independent Gaussian standard deviation with unbounded
samples clipped by the environment adapter; reward scaling by the running
standard deviation of the return instead of MAPPO's value normalisation
(PopArt or ValueNorm), both of which keep value targets of order one; and the
concatenation of the local observations as MAPPO's global state, not an
environment-specific state.

\begin{table}[htbp]
  \centering\small
  \caption{Fields of \code{PPOConfig}, in addition to those of
    \code{OnPolicyConfig} and \code{AlgorithmConfig}
    (Table~\ref{tab:marl-config-common}).}
  \label{tab:marl-ppo-config}
  \begin{tabularx}{\textwidth}{@{}>{\raggedright\arraybackslash}p{3.7cm}>{\raggedright\arraybackslash}p{1.5cm}>{\raggedright\arraybackslash}X@{}}
    \toprule
    Field & Default & Meaning \\
    \midrule
    \code{n\_epochs} & 10 & passes over each rollout \\
    \code{num\_minibatches} & 4 & minibatches per epoch (a sample is one copy at one step with all its agents) \\
    \code{clip\_range} & 0.2 & ratio clipping $\epsilon$ \\
    \code{value\_clip\_range} & 0.2 & value clipping $\epsilon_V$ around the pre-update value; \code{None} disables it \\
    \code{huber\_delta} & \code{None} & Huber threshold $\eta$ of the critic loss; \code{None} uses $e^2/2$ \\
    \code{ent\_coef} & 0.0 & entropy weight $c_H$ \\
    \code{normalize\_advantages} & \code{True} & standardise the advantages over the rollout \\
    \code{critic\_lr} & \code{None} & critic learning rate; \code{None} uses \code{lr} \\
    \code{critic\_hidden\_sizes} & \code{None} & hidden widths of the critic; \code{None} uses \code{hidden\_sizes} \\
    \code{anneal\_lr} & \code{False} & decay both learning rates linearly to zero \\
    \code{target\_kl} & \code{None} & stop the epochs once the mean approximate KL of an epoch exceeds $1.5\,\times$ this value \\
    \code{log\_std\_init} & $-0.5$ & initial log standard deviation of the Gaussian policy, in units of half the action range \\
    \code{reward\_source} & \code{None} & \code{"team"} or \code{"agent"}; \code{None} selects \code{"agent"} for IPPO and \code{"team"} for MAPPO \\
    \bottomrule
  \end{tabularx}
\end{table}

% ---------------------------------------------------------------------------
\section{MADDPG and MATD3}\label{sec:marl-ddpg}

MADDPG (Lowe et al., 2017) extends DDPG (Lillicrap et al., 2016) to several
agents: decentralised deterministic actors $\mu_i(o_i)$ are trained with the
deterministic policy gradient (Silver et al., 2014) of centralised critics
$Q_i(s, a_1, \dots, a_n)$ that see the global state and the joint action.
MATD3 (Ackermann et al., 2019) carries the three corrections of TD3
(Fujimoto et al., 2018) over to the centralised critics. Both learn from a
replay buffer and support continuous actions only.

\paragraph{Normalised action units.} The actors output actions in $[-1, 1]$
through a $\tanh$; the per-agent bounds of the environment map them to
physical units, $a_i = c_i + h_i \odot \tilde a_i$, and actions from the
replay buffer are mapped back, $\tilde a_i = (a_i - c_i) \oslash h_i$. Critics,
exploration noise and target smoothing all work in these normalised units, so
one set of hyperparameters suits environments with very different action
ranges.

\paragraph{Critic inputs.} With shared parameters (the default) one actor and
one critic serve all agents, with the one-hot identifier $e_i$ appended to the
actor input $x_i = [\hat o_i, e_i]$ and to the critic input; the shared critic
evaluated at agent $i$'s identifier is $Q_i$. With \code{share\_parameters=False}
every agent has its own actor and critic. By default
(\code{critic\_local\_inputs=True}) the input of $Q_i$ also repeats agent $i$'s
own observation and action,
\begin{equation}
  z_i(s, \tilde a) = \bigl[\,s,\; \tilde a_1, \dots, \tilde a_n,\; \hat o_i,\;
    \tilde a_i,\; e_i\,\bigr], \qquad Q_i(s, \tilde a) = Q_\phi\bigl(z_i(s, \tilde a)\bigr).
  \label{eq:marl-critic-input}
\end{equation}
This is a documented deviation from the original critic input
$[s, a_1, \dots, a_n]$. The input is still a function of $(s, a)$ and $i$ only
(the ``agent-specific global state'' of Yu et al., 2022), but a shared critic
no longer has to learn from the identifier which of the $n$ blocks of $s$ and
$a$ belong to agent $i$. According to the implementation notes, MADDPG does not
learn on \envid{Consensus-v0} (8 agents) without these inputs and approaches
the Laplacian baseline with them. \code{critic\_local\_inputs=False} restores
the original input.

\paragraph{Critic update.} On a replay batch of $N$ transitions
$(o, a, r, o', d)$, with $d$ the \code{terminated} flag (truncated episodes
bootstrap from their stored final observations) and $r_i$ agent $i$'s own
reward (\code{reward\_source="agent"}, the default, as in MADDPG) or the team
reward (\code{"team"}), the target actions and the TD targets are
\begin{align}
  \tilde a'_j &= \mu_{\theta'}(x'_j) &&\text{(MADDPG)}, \nonumber\\
  \tilde a'_j &= \operatorname{clip}\bigl(\mu_{\theta'}(x'_j)
      + \operatorname{clip}(\varepsilon_j, -c, c),\, -1, 1\bigr),
      \quad \varepsilon_j \sim \mathcal{N}(0, \sigma^2 I) &&\text{(MATD3)},
      \label{eq:marl-target-actions}\\
  y_i &= \kappa\, r_i + \gamma\,(1 - d)\, \min_{k} Q'_{i,k}(s', \tilde a'), \nonumber
\end{align}
with the target actor $\mu_{\theta'}$, the target critics $Q'_{i,k}$ ($k = 1$
for MADDPG, $k \in \{1, 2\}$ for MATD3), $\kappa$ = \code{reward\_scale},
$\sigma$ = \code{target\_noise} and $c$ = \code{target\_noise\_clip}. The
critics minimise
\begin{equation}
  L^{Q}(\phi) = \sum_k \frac{1}{N n} \sum_{b, i}
      \ell\bigl(Q_{i,k}(s_b, \tilde a_b) - y_{b,i}\bigr),
  \label{eq:marl-ddpg-critic}
\end{equation}
with $\ell(e) = e^2$ (\code{critic\_loss="mse"}, as in the papers) or, with
\code{critic\_loss="huber"}, $\ell(e) = e^2$ for $|e| \le 1$ and $2|e| - 1$
beyond, which bounds the gradient of rare, very large errors such as the trip
penalty of \envid{PowerGrid-v0}. The gradient norm is clipped at
\code{max\_grad\_norm}.

\paragraph{Actor update.} Actor $i$ follows the deterministic policy gradient
of its (first) critic, with its own action replaced by $\mu_\theta(x_i)$ and
the other agents' actions taken from the replay batch:
\begin{equation}
  L^{\mu}(\theta) = -\frac{1}{N n} \sum_{b, i}
      Q_{i,1}\bigl(s_b,\; (\tilde a_{b,1}, \dots, \mu_\theta(x_{b,i}), \dots,
      \tilde a_{b,n})\bigr).
  \label{eq:marl-ddpg-actor}
\end{equation}
The replacement applies to the joint action and, with local inputs, to the
own-action slot of \eqref{eq:marl-critic-input}. All agents are updated in one
batched pass (row $i$ of the batched critic input carries agent $i$'s
replacement); the critic parameters are frozen during this step, so only the
actor changes. After every actor update the target networks follow by Polyak
averaging, $\theta' \leftarrow (1 - \tau)\theta' + \tau\theta$ and
$\phi' \leftarrow (1 - \tau)\phi' + \tau\phi$. MADDPG updates the actor after
every critic update; MATD3 updates the actor and the targets only every
\code{policy\_delay} critic updates, which together with the twin critics and
the target smoothing of \eqref{eq:marl-target-actions} are the three TD3
corrections.

\paragraph{Exploration.} During training, $\tilde a_i =
\operatorname{clip}(\mu_\theta(x_i) + \sigma_t \xi_i, -1, 1)$ with $\xi_i \sim
\mathcal{N}(0, I)$, drawn from the algorithm's generator. The noise scale
$\sigma_t$ is \code{exploration\_noise} or, when
\code{final\_exploration\_noise} is set, decays linearly with the environment
steps $t$ (summed over copies),
\begin{equation}
  \sigma_t = \sigma_0 + \min(1, t / T_\sigma)\,(\sigma_1 - \sigma_0),
\end{equation}
where $\sigma_0$ is \code{exploration\_noise}, $\sigma_1$ is
\code{final\_exploration\_noise} and $T_\sigma$ is the length
\code{noise\_decay\_steps} of the decay; the property
\code{exploration\_scale} returns the current value. The
warm-up actions are uniform in the bounds.

\paragraph{Observation normalisation.} Optional
(\code{normalize\_observations}): the statistics are computed once, from all
observations in the replay buffer at the first gradient step (the warm-up
data), and frozen afterwards, so that stored transitions and learned values
stay consistent. Every update records \code{critic\_loss}, \code{actor\_loss}
(of the last actor update), \code{q\_mean}, \code{target\_mean} and
\code{noise\_scale}.

\begin{table}[htbp]
  \centering\small
  \caption{Fields of \code{DDPGConfig}, in addition to those of
    \code{OffPolicyConfig} and \code{AlgorithmConfig}
    (Table~\ref{tab:marl-config-common}). Noise scales are in units of half the
    action range.}
  \label{tab:marl-ddpg-config}
  \begin{tabularx}{\textwidth}{@{}>{\raggedright\arraybackslash}p{4.3cm}>{\raggedright\arraybackslash}p{1.5cm}>{\raggedright\arraybackslash}X@{}}
    \toprule
    Field & Default & Meaning \\
    \midrule
    \code{actor\_lr}, \code{critic\_lr} & $10^{-3}$ & Adam learning rates \\
    \code{critic\_hidden\_sizes} & \code{None} & hidden widths of the critics; \code{None} uses \code{hidden\_sizes} (which always sets the actors) \\
    \code{exploration\_noise} & 0.1 & standard deviation $\sigma_0$ of the exploration noise \\
    \code{final\_exploration\_noise} & \code{None} & when set, the noise decays linearly to this value $\sigma_1$ \\
    \code{noise\_decay\_steps} & 0 & length $T_\sigma$ of the decay in environment steps \\
    \code{reward\_source} & \code{"agent"} & \code{"agent"}: critic $i$ learns from agent $i$'s reward; \code{"team"}: from the team reward \\
    \code{normalize\_observations} & \code{False} & standardise observations with statistics of the warm-up data, frozen at the first gradient step \\
    \code{critic\_local\_inputs} & \code{True} & append agent $i$'s own observation and action to the input of $Q_i$ \eqref{eq:marl-critic-input} \\
    \code{critic\_loss} & \code{"mse"} & \code{"mse"} or \code{"huber"} \\
    \code{policy\_delay} & 2 & MATD3 only: critic updates per actor and target update \\
    \code{target\_noise} & 0.2 & MATD3 only: standard deviation $\sigma$ of the target smoothing noise \\
    \code{target\_noise\_clip} & 0.5 & MATD3 only: clipping bound $c$ of the smoothing noise \\
    \bottomrule
  \end{tabularx}
\end{table}

% ---------------------------------------------------------------------------
\section{IQL, VDN and QMIX}\label{sec:marl-value}

The value-based methods learn per-agent \emph{utilities}
$Q_\theta(x_i) \in \mathbb{R}^{|\mathcal{A}|}$, one value per action of agent
$i$, with the DQN machinery (Mnih et al., 2015): a replay buffer, target
networks and $\epsilon$-greedy exploration. Every agent acts greedily on its
own utility, $a_i = \arg\max_a Q_\theta(x_i)[a]$, from its own observation. The
three methods differ only in how the utilities are trained. With shared
parameters (default) one utility network serves all agents, with the one-hot
identifier appended to the observation as in the QMIX paper;
\code{share\_parameters=False} gives one network per agent. The observations
enter the utility network unnormalised.

On a replay batch $(o, a, r, o', d)$ the chosen utilities are
$q_i = Q_\theta(x_i)[a_i]$. The greedy next actions are selected by the online
network and evaluated by the target network $\theta^-$ (double Q-learning, van
Hasselt et al., 2016, \code{double\_q=True}),
\begin{equation}
  \hat a'_i = \arg\max_a Q_\theta(x'_i)[a], \qquad
  q'_i = Q_{\theta^-}(x'_i)[\hat a'_i],
\end{equation}
or both by the target network with \code{double\_q=False}. With $\kappa$ =
\code{reward\_scale}, $d$ the \code{terminated} flag and $\ell$ the Huber loss
(smooth L1 with threshold 1, $\ell(e) = e^2/2$ for $|e| < 1$ and $|e| - 1/2$
otherwise; \code{loss="huber"}) or the squared error $\ell(e) = e^2$
(\code{loss="mse"}):

\begin{itemize}
  \item \textbf{IQL} (Tan, 1993, here with DQN agents): every agent regresses
    its utility on its own reward $r_i$ with its own target; there is no
    mixing and no centralised information,
    \begin{equation}
      y_i = \kappa\, r_i + \gamma (1 - d)\, q'_i, \qquad
      L = \frac{1}{N n} \sum_{b, i} \ell(q_{b,i} - y_{b,i}).
    \end{equation}
  \item \textbf{VDN} (Sunehag et al., 2018): the team value is the sum of the
    utilities, trained end to end on the team reward $r$ with one target,
    \begin{equation}
      Q_\text{tot} = \sum_i q_i, \qquad
      y = \kappa\, r + \gamma (1 - d) \sum_i q'_i, \qquad
      L = \frac{1}{N} \sum_b \ell(Q_{\text{tot},b} - y_b).
    \end{equation}
  \item \textbf{QMIX} (Rashid et al., 2018): the team value is a monotonic
    function of the utilities, computed by a mixing network whose weights are
    generated from the global state $s$ by hypernetworks,
    \begin{align}
      Q_\text{tot}(q; s) &= w_2(s)^{\top} \operatorname{elu}\bigl(W_1(s)^{\top} q
        + b_1(s)\bigr) + b_2(s), \label{eq:marl-qmix}\\
      W_1(s) &= \bigl|h_{W_1}(s)\bigr| \in \mathbb{R}^{n \times m}, \quad
      w_2(s) = \bigl|h_{w_2}(s)\bigr| \in \mathbb{R}^{m}, \nonumber\\
      y &= \kappa\, r + \gamma (1 - d)\, Q^-_\text{tot}(q'; s'), \qquad
      L = \frac{1}{N} \sum_b \ell\bigl(Q_\text{tot}(q_b; s_b) - y_b\bigr), \nonumber
    \end{align}
    with $m$ = \code{mixer\_embed\_dim}. The hypernetworks $h_{W_1}$ and
    $h_{w_2}$ are MLPs with one hidden layer of \code{hypernet\_hidden} ReLU
    units, $b_1$ is a linear function of $s$ and $b_2$ an MLP with one hidden
    layer of $m$ ReLU units; the absolute values make the weights
    non-negative, and since $\operatorname{elu}$ is increasing,
    $\partial Q_\text{tot} / \partial q_i \ge 0$. The target mixer $Q^-_\text{tot}$
    is a Polyak-averaged or periodically copied twin of the mixer.
\end{itemize}

Because the VDN and QMIX mixing is monotonic in every utility, the greedy joint
action $\arg\max_{a} Q_\text{tot}$ is obtained by every agent maximising its own
utility (the individual-global-max property). The centralised target is
therefore computed from per-agent maxima, the mixers are used only in
training, and execution stays decentralised. The global state of QMIX is the
concatenation of all agents' observations (Section~\ref{sec:marl-spec}).

\paragraph{Optimisation and exploration.} One Adam optimiser (learning rate
\code{lr}) updates the utility network and the mixer, with the gradient norm of
all their parameters clipped at \code{max\_grad\_norm}. After every gradient
step the target networks follow by Polyak averaging with \code{tau}
(\code{target\_update\_interval=0}, the default), or the online weights are
copied into them every $K$ gradient steps (\code{target\_update\_interval}
$= K > 0$, the original DQN scheme). While training, every agent of every copy
independently takes a uniformly random action with probability
\begin{equation}
  \epsilon(t) = \epsilon_\text{start} + \min(1, t / T_\epsilon)\,
      (\epsilon_\text{end} - \epsilon_\text{start}),
\end{equation}
where $t$ counts environment steps (summed over copies) from the start of
training, including the warm-up, and $T_\epsilon$ =
\code{epsilon\_decay\_steps} ($0$ uses $\epsilon_\text{end}$ from the start);
\code{epsilon\_at(t)} and the property \code{epsilon} return the value. Every
update records \code{td\_loss}, \code{q\_mean} (the trained value,
$Q_\text{tot}$ or $q_i$, in units of the scaled reward), \code{target\_mean},
\code{grad\_norm} (before clipping) and \code{epsilon}. The methods
\code{q\_values(obs, *, target=False)}, \code{mix(agent\_q, obs, *,
target=False)}, \code{td\_target(batch)} and \code{td\_loss(batch)} expose the
pieces of the update.

\paragraph{Stated simplifications.} The VDN and QMIX papers use recurrent
(LSTM or GRU) agent networks trained on whole episodes sampled from an
episode replay, to cope with partial observability. Here the agent networks
are feed-forward and trained on single transitions from a transition replay:
the observations of the discrete environments (the local message and queue
windows of LineMsg and WirelessComm, the piston and ball features of
Pistonball) are Markov enough for the demonstrations, and the feed-forward
version is much cheaper on a CPU. Further standard choices: one-step TD
targets, double Q-learning by default, Adam instead of RMSprop, and
Polyak-averaged target networks by default. IQL with environments that report
no per-agent rewards gives every agent the team reward, and with a single
agent it is (double) DQN.

\begin{table}[htbp]
  \centering\small
  \caption{Fields of \code{QLearningConfig} (IQL, VDN, QMIX), in addition to
    those of \code{OffPolicyConfig} and \code{AlgorithmConfig}
    (Table~\ref{tab:marl-config-common}).}
  \label{tab:marl-value-config}
  \begin{tabularx}{\textwidth}{@{}>{\raggedright\arraybackslash}p{4.3cm}>{\raggedright\arraybackslash}p{1.5cm}>{\raggedright\arraybackslash}X@{}}
    \toprule
    Field & Default & Meaning \\
    \midrule
    \code{lr} & $5 \cdot 10^{-4}$ & Adam learning rate of the utility network and the mixer \\
    \code{epsilon\_start}, \code{epsilon\_end} & 1.0, 0.05 & exploration probability at the start and after the decay \\
    \code{epsilon\_decay\_steps} & 50\,000 & environment steps $T_\epsilon$ of the linear decay \\
    \code{double\_q} & \code{True} & select the next action with the online network, evaluate it with the target network \\
    \code{target\_update\_interval} & 0 & 0: Polyak averaging with \code{tau} after every gradient step; $K > 0$: copy the online weights every $K$ gradient steps \\
    \code{loss} & \code{"huber"} & TD loss, \code{"huber"} (threshold 1) or \code{"mse"} \\
    \code{mixer\_embed\_dim} & 32 & QMIX: width $m$ of the mixing layer \\
    \code{hypernet\_hidden} & 64 & QMIX: hidden width of the weight hypernetworks \\
    \bottomrule
  \end{tabularx}
\end{table}

% ---------------------------------------------------------------------------
\section{Presets and results}\label{sec:marl-results}

The package \pkg{marl\_algorithms.presets} holds a table
\code{PRESETS[algorithm][env\_id]} of tuned settings, collected from one module
per family (\pkg{presets.ppo}, \pkg{presets.ddpg} and
\pkg{presets.q\_learning}, each with the notes from tuning). A preset fixes the environment steps (\code{"total\_steps"}),
the number of copies (\code{"num\_envs"}), configuration overrides
(\code{"config"}) and environment arguments (\code{"env\_kwargs"}, empty for
all current presets). The presets were tuned to train in at most about two
minutes on one CPU core and to show clear learning.
Table~\ref{tab:marl-presets} lists all 17 of them, as printed by
\code{marl-train presets}, with their configuration overrides.

\begin{table}[htbp]
  \centering\small
  \caption{Presets (\code{list\_presets()}); ``Steps'' is the preset's
    \code{total\_steps}, ``Copies'' its \code{num\_envs}.}
  \label{tab:marl-presets}
  \begin{tabularx}{\textwidth}{@{}l>{\raggedright\arraybackslash}p{2.9cm}rr>{\raggedright\arraybackslash}X@{}}
    \toprule
    Algorithm & Environment & Steps & Copies & Configuration overrides \\
    \midrule
    IPPO  & \envid{PowerGrid-v0} & 200\,000 & 32 & \code{rollout\_length=64}, \code{n\_epochs=5}, \code{lr=1e-3}, \code{gamma=0.95}, \code{log\_std\_init=-1.5} \\
    IPPO  & \envid{Platoon-v0}   & 450\,000 & 32 & \code{rollout\_length=64}, \code{n\_epochs=5}, \code{lr=3e-4}, \code{gamma=0.99}, \code{log\_std\_init=-1.0}, \code{normalize\_observations=False}, \code{value\_clip\_range=None} \\
    IPPO  & \envid{Consensus-v0} & 400\,000 & 32 & \code{rollout\_length=64}, \code{n\_epochs=5}, \code{lr=1e-3}, \code{gamma=0.95}, \code{normalize\_observations=False} \\
    IPPO  & \envid{LineMsg-v0}   & 100\,000 & 32 & \code{rollout\_length=50}, \code{n\_epochs=5}, \code{lr=1e-3}, \code{ent\_coef=0.01} \\
    MAPPO & \envid{PowerGrid-v0} & 250\,000 & 32 & as IPPO on \envid{PowerGrid-v0} \\
    MAPPO & \envid{Platoon-v0}   & 450\,000 & 32 & as IPPO on \envid{Platoon-v0}, and \code{reward\_source="agent"} \\
    MAPPO & \envid{Consensus-v0} & 400\,000 & 32 & \code{rollout\_length=64}, \code{n\_epochs=5}, \code{lr=1e-3}, \code{gamma=0.9}, \code{reward\_source="agent"}, \code{normalize\_observations=False} \\
    MAPPO & \envid{LineMsg-v0}   & 100\,000 & 32 & as IPPO on \envid{LineMsg-v0} \\
    MADDPG & \envid{PowerGrid-v0} & 64\,000 & 16 & \code{critic\_hidden\_sizes=(128, 128)}, \code{batch\_size=64}, \code{warmup\_steps=3200}, \code{reward\_source="team"}, \code{critic\_loss="huber"} \\
    MADDPG & \envid{Consensus-v0} & 96\,000 & 16 & \code{critic\_hidden\_sizes=(128, 128)}, \code{batch\_size=64}, \code{warmup\_steps=3200}, \code{gamma=0.95}, \code{reward\_scale=0.01}, \code{normalize\_observations=True}, \code{exploration\_noise=0.2}, \code{final\_exploration\_noise=0.02}, \code{noise\_decay\_steps=96000} \\
    MATD3 & \envid{PowerGrid-v0} & 64\,000 & 16 & as MADDPG, and \code{exploration\_noise=0.3}, \code{final\_exploration\_noise=0.05}, \code{noise\_decay\_steps=64000} \\
    MATD3 & \envid{Consensus-v0} & 96\,000 & 16 & as MADDPG on \envid{Consensus-v0} \\
    IQL, VDN, QMIX & \envid{LineMsg-v0} & 25\,000 & 16 & \code{batch\_size=128}, \code{lr=1e-3}, \code{epsilon\_decay\_steps=12500} \\
    VDN  & \envid{WirelessComm-v1} & 120\,000 & 16 & \code{batch\_size=128}, \code{lr=1e-3}, \code{gamma=0.8}, \code{epsilon\_decay\_steps=30000} \\
    QMIX & \envid{WirelessComm-v1} & 100\,000 & 16 & \code{batch\_size=128}, \code{lr=1e-3}, \code{gamma=0.8}, \code{epsilon\_decay\_steps=45000} \\
    \bottomrule
  \end{tabularx}
\end{table}

The source code documents the reasons for these settings. On
\envid{PowerGrid-v0} the PPO policies need little initial exploration noise
(\code{log\_std\_init=-1.5}): with more noise the mean policy learns to
compensate its own previous actions, which it observes, and performs poorly
when executed without noise. On \envid{Platoon-v0} and \envid{Consensus-v0} the
critics learn each agent's own reward, which isolates the effect of an agent's
own command, and the observations are used unnormalised, because they are
already of order one (Platoon) or in arena units (Consensus) and running
statistics dominated by the large errors of early training would hide the
small errors that matter once the task is nearly solved. MADDPG and MATD3 use
the team reward on \envid{PowerGrid-v0}, where each bus sees only about
1/16 of the benefit of its frequency support but pays its full control cost
under the per-agent rewards, and the Huber loss keeps the rare trip
transitions from dominating the critic updates. On \envid{WirelessComm-v1}
packets are delivered in the step they are sent, so a short horizon
($\gamma = 0.8$) suffices. IQL has no \envid{WirelessComm-v1} preset, for the
reason discussed below.

\paragraph{Results.} Table~\ref{tab:marl-results-table} reports the results
of \code{benchmarks/benchmark\_marl.py}, which applies one protocol to every
preset: train with \code{train\_preset(algorithm, env\_id, seed=0)} on one CPU
core (\code{OMP\_NUM\_THREADS=1}, one PyTorch thread), then evaluate uniformly
random actions, the trained deterministic policy and the baseline controller
of the environment with \code{env\_lib.evaluate} on the same 64 episodes (seed
1) of a 64-copy vector environment created with \code{make\_vector\_env}. The
script appends one JSON record per preset to the file given by \code{-{}-out}
(default \code{renders/marl\_benchmark.jsonl}); \code{-{}-only ALGO:ENV ...}
selects presets, \code{-{}-shard K/N} runs every $N$-th preset starting at
$K$, so that shards can run in parallel processes (for example pinned to
different cores with \code{taskset}), and \code{-{}-report FILE} prints the
table of a result file, as Markdown with \code{-{}-markdown}.

\begin{shellcode}
python benchmarks/benchmark_marl.py --out marl.jsonl          # all presets
python benchmarks/benchmark_marl.py --only mappo:PowerGrid-v0 qmix:LineMsg-v0
taskset -c 0 python benchmarks/benchmark_marl.py --shard 0/2 --out marl.jsonl &
taskset -c 1 python benchmarks/benchmark_marl.py --shard 1/2 --out marl.jsonl &
python benchmarks/benchmark_marl.py --report marl.jsonl --markdown
\end{shellcode}

\begin{table}[htbp]
  \centering\small
  \caption{Results of \code{benchmarks/benchmark\_marl.py}: mean team return of
    64 evaluation episodes (seed 1, 64-copy vector environment) after training
    with the preset and seed 0; higher is better. ``CPU'' is the processor
    time of training on one core. ``Env steps'' are the steps actually
    collected, a multiple of the steps per rollout (on-policy) or of the
    number of copies (off-policy). The baselines are droop control
    (PowerGrid), cooperative adaptive cruise control (Platoon), the Laplacian
    protocol (Consensus), always relay (LineMsg) and the collision-free
    schedule (WirelessComm).}
  \label{tab:marl-results-table}
  \begin{tabular}{@{}llrrrrr@{}}
    \toprule
    Environment & Algorithm & Env steps & CPU [s] & Random & Trained & Baseline \\
    \midrule
    \envid{PowerGrid-v0} & IPPO   & 200\,704 & 69  & $-629.3$ & $-2.37$ & $-0.79$ \\
                         & MAPPO  & 251\,904 & 106 & $-629.3$ & $-0.89$ & $-0.79$ \\
                         & MADDPG & 64\,000  & 70  & $-629.3$ & $-5.69$ & $-0.79$ \\
                         & MATD3  & 64\,000  & 87  & $-629.3$ & $-5.63$ & $-0.79$ \\
    \midrule
    \envid{Platoon-v0}   & IPPO   & 450\,560 & 88  & $-5350$ & $-84.8$ & $-46.4$ \\
                         & MAPPO  & 450\,560 & 102 & $-5350$ & $-65.0$ & $-46.4$ \\
    \midrule
    \envid{Consensus-v0} & IPPO   & 401\,408 & 83  & $-18041$ & $-1550$ & $-1518$ \\
                         & MAPPO  & 401\,408 & 96  & $-18041$ & $-1586$ & $-1518$ \\
                         & MADDPG & 96\,000  & 72  & $-18041$ & $-1681$ & $-1518$ \\
                         & MATD3  & 96\,000  & 83  & $-18041$ & $-1737$ & $-1518$ \\
    \midrule
    \envid{LineMsg-v0}   & IPPO   & 100\,800 & 26 & 43.2 & 95.0 & 95.0 \\
                         & MAPPO  & 100\,800 & 26 & 43.2 & 95.0 & 95.0 \\
                         & IQL    & 25\,008  & 9  & 43.2 & 95.0 & 95.0 \\
                         & VDN    & 25\,008  & 10 & 43.2 & 95.0 & 95.0 \\
                         & QMIX   & 25\,008  & 16 & 43.2 & 95.0 & 95.0 \\
    \midrule
    \envid{WirelessComm-v1} & VDN  & 120\,000 & 74  & 139.2 & 338.0 & 339.6 \\
                            & QMIX & 100\,000 & 101 & 139.2 & 319.1 & 339.6 \\
    \bottomrule
  \end{tabular}
\end{table}

\paragraph{Discussion.} Every preset learns: each trained policy is far better
than random actions, and several match their baseline.

\begin{itemize}
  \item \emph{LineMsg.} All five discrete-action methods learn the optimal
    ``always relay'' policy (return 95.0, the same as the baseline); the
    value-based methods need a quarter of the environment steps of the PPO
    methods and 9 to 16 seconds.
  \item \emph{Consensus.} IPPO and MAPPO come within about 2 and 5\,\% of the
    Laplacian protocol, MADDPG and MATD3, with a quarter of the environment
    steps, within about 11 and 14\,\%. The Laplacian protocol is a strong reference
    here: it is the classical solution of the consensus problem on a
    connected graph.
  \item \emph{PowerGrid.} MAPPO reaches $-0.89$ against $-0.79$ for droop
    control, IPPO $-2.37$. MADDPG and MATD3 improve the return from $-629$ to
    about $-5.7$ but remain far from droop control. Fitting the actions of the
    trained deterministic policies to the local frequency deviation $\omega_i$
    (measured for this chapter over 200 steps on 16 copies) gives slopes of
    $-0.09$ (MADDPG)
    and $-0.14$ (MATD3), against the droop gain of $-0.25$, with correlations
    of only $-0.15$ and $-0.12$: the actors have learned a partial droop
    feedback, weaker than droop control, combined with a sizeable part of the
    action that does not follow the local frequency deviation and costs
    control effort.
  \item \emph{Platoon.} IPPO and MAPPO reduce the cost from about 5\,350 to 85
    and 65, which is still 83 and 40\,\% above the cost of cooperative
    adaptive cruise control. According to the tuning notes in the source,
    MAPPO with the team reward and normalised observations does not learn this
    task (every episode ends in a collision); the per-agent rewards and
    unnormalised observations of the preset are essential.
  \item \emph{WirelessComm.} VDN comes within 0.5\,\% of the collision-free
    schedule. IQL has no preset here: trained with the VDN settings
    (120\,000 steps), it reached only 152.3 on the same evaluation episodes,
    hardly better than random actions. This is a structural failure rather
    than a matter of tuning. The per-agent reward is 1 for a delivered packet
    and 0 otherwise, so a collision costs the transmitting agent nothing
    compared with staying idle; for an independent learner transmitting is
    never worse than waiting, the harm it does to the other agent's delivery is
    invisible in its own reward, and the agents keep transmitting and
    colliding. VDN and QMIX learn from the team reward, in which a collision
    costs a delivery, and learn to leave a contested access point to one agent.
    The QMIX result depends on the seed. With seed 0 it converges to a fixed
    assignment in which each of the nine access points is used by exactly one
    agent, the one to its upper right, and the seven agents of the first column
    and the last row never deliver a packet (6 collisions in the 64 evaluation
    episodes); the baseline additionally lets the border agents borrow idle
    access points. With the same preset and seeds 1 and 2, QMIX reached 344.1
    and 332.7 on the same evaluation episodes (measured for this chapter),
    with 14 and 12 agents delivering packets.
\end{itemize}

The off-policy methods need three to four times fewer environment steps than the
PPO methods on the continuous tasks, but spend more computation per step, so
the training times are similar. The results come from single training runs
(seed 0); a comparison of methods for a publication needs several seeds per
method and confidence intervals, for example with
\code{plot\_learning\_curves} of \pkg{plotkit}
(Section~\ref{sec:plotkit-curves}).

% ---------------------------------------------------------------------------
\section{The demonstration script}\label{sec:marl-demo}

\code{examples/algorithms/marl\_training\_demo.py} trains one algorithm of every family
with its preset, each on a different environment:

\begin{itemize}
  \item MAPPO on \envid{PowerGrid-v0}: continuous actions, 16 agents, on-policy;
  \item MADDPG on \envid{Consensus-v0}: continuous actions, 8 agents, off-policy;
  \item QMIX on \envid{LineMsg-v0}: discrete actions, 10 agents, value-based;
  \item VDN on \envid{WirelessComm-v1}: discrete actions, 16 agents, value-based.
\end{itemize}

Each run trains on batched copies created with \code{make\_vector\_env}
(native batches for PowerGrid and Consensus), then evaluates random actions,
the trained policy and the baseline controller on the same seeded episodes (a
vector environment of \code{min(episodes, 64)} copies, seed \code{seed + 1}).
The script
prints one line per run and a results table, and saves the learning curves
(training returns over environment steps, with the random and baseline
evaluation returns as reference lines) as a PNG file. The full suite takes a
few minutes on one CPU core.

\begin{table}[htbp]
  \centering\small
  \caption{Options of \code{examples/algorithms/marl\_training\_demo.py}.}
  \label{tab:marl-demo-options}
  \begin{tabularx}{\textwidth}{@{}>{\raggedright\arraybackslash}p{2.9cm}>{\raggedright\arraybackslash}p{4.6cm}>{\raggedright\arraybackslash}X@{}}
    \toprule
    Option & Default & Meaning \\
    \midrule
    \code{-{}-algo}, \code{-{}-env} & the suite & train only this preset (both are required together) \\
    \code{-{}-quick} & off & smoke test: 2\,000 steps (PPO) or 1\,500 steps with a 500-step warm-up (off-policy) on 8 copies \\
    \code{-{}-seed} & 0 & training seed; the evaluation uses \code{seed + 1} \\
    \code{-{}-episodes} & 64 & evaluation episodes per policy \\
    \code{-{}-plot} & \code{renders/marl\_training.png} & learning-curve figure; an empty string skips it \\
    \code{-{}-gif} & none & record the trained policy of a single run (\code{-{}-algo}/\code{-{}-env}) \\
    \code{-{}-threads} & 1 & PyTorch threads \\
    \bottomrule
  \end{tabularx}
\end{table}

\begin{shellcode}
python examples/algorithms/marl_training_demo.py                 # the suite, a few minutes
python examples/algorithms/marl_training_demo.py --algo mappo --env PowerGrid-v0
python examples/algorithms/marl_training_demo.py --algo qmix --env LineMsg-v0 \
    --gif renders/qmix.gif
python examples/algorithms/marl_training_demo.py --quick --plot ''   # smoke test
\end{shellcode}

Table~\ref{tab:marl-demo-results} shows the results of the full suite with its
default arguments and Figure~\ref{fig:marl-training} the learning curves it
saves. The evaluation protocol is that of
Table~\ref{tab:marl-results-table}, and the returns agree with it to the digits
shown there.

\begin{table}[htbp]
  \centering\small
  \caption{Output of \code{examples/algorithms/marl\_training\_demo.py} with default
    arguments: training seed 0 on one CPU core (``Time'' is the wall-clock
    training time), 64 evaluation episodes with seed 1 on a 64-copy vector
    environment.}
  \label{tab:marl-demo-results}
  \begin{tabular}{@{}llrrrrr@{}}
    \toprule
    Algorithm & Environment & Env steps & Time [s] & Random & Trained & Baseline \\
    \midrule
    MAPPO  & \envid{PowerGrid-v0}    & 251\,904 & 106.6 & $-629.3$  & $-0.894$ & $-0.786$ \\
    MADDPG & \envid{Consensus-v0}    & 96\,000  & 76.1  & $-18041$  & $-1681$  & $-1518$ \\
    QMIX   & \envid{LineMsg-v0}      & 25\,008  & 16.6  & 43.2      & 95.0     & 95.0 \\
    VDN    & \envid{WirelessComm-v1} & 120\,000 & 78.8  & 139.2     & 338.0    & 339.6 \\
    \bottomrule
  \end{tabular}
\end{table}

\begin{figure}[htbp]
  \centering
  \includegraphics[width=\textwidth]{marl_training.png}
  \caption{Learning curves saved by \code{examples/algorithms/marl\_training\_demo.py}
    (default suite, seed 0): mean return of the exploratory training episodes
    over environment steps, with the evaluation returns of random actions
    (dotted) and of the baseline controller (dashed) for reference. The
    training returns include the random warm-up and the exploration (policy
    noise or $\epsilon$-greedy actions) and are averaged over the few episodes
    that finish in each interval, so they are not directly comparable with
    the evaluation returns of the deterministic policy in
    Table~\ref{tab:marl-demo-results}.}
  \label{fig:marl-training}
\end{figure}

% ---------------------------------------------------------------------------
\section{Adding an algorithm}\label{sec:marl-extending}

A new algorithm reuses the core and implements only what is specific to it.
Inside the package it gets its own module in the folder of its family (or a
new family folder with a \code{common.py} when it shares code with a sibling):

\begin{enumerate}
  \item Define a configuration dataclass that derives from
    \code{OnPolicyConfig} or \code{OffPolicyConfig} (validate new fields in
    \code{\_\_post\_init\_\_}, after calling the parent's).
  \item Subclass \code{OnPolicyAlgorithm} or \code{OffPolicyAlgorithm} and set
    the class attributes \code{name}, \code{action\_kinds} and
    \code{config\_class}.
  \item Implement the hooks: \code{\_build()} (networks and optimisers; it runs
    under the seeded, forked PyTorch generator), \code{act(obs,
    deterministic=False)}, \code{\_modules()} (what \code{save} stores) and,
    for an on-policy method, \code{\_rollout\_step(obs)} and
    \code{\_update(rollout)}, optionally \code{\_rollout\_fields()} and
    \code{\_observe\_transition(transition)}; for an off-policy method,
    \code{\_explore(obs)} and \code{\_update(batch)}. Extra state such as
    normalisers goes into \code{\_extra\_state()} and
    \code{\_load\_extra\_state(state)}. Use \code{self.np\_rng} and
    \code{self.torch\_rng} for all randomness, \code{self.agent\_inputs} for
    per-agent network inputs and \code{PerAgent} for shared or separate
    networks, and batch over copies and agents.
  \item Register it: add an \code{AlgorithmInfo} record (name,
    \code{"module:Class"} entry point, family, action kinds, summary and
    reference) to the tuple in \pkg{marl\_algorithms.registry}, and optionally a
    lazy export in \code{marl\_algorithms/\_\_init\_\_.py}. \code{train},
    \code{make\_algorithm}, \code{Algorithm.load} and \code{marl-train} then
    find it. An unregistered class can still be trained directly
    (\code{MyAlgorithm(envs, seed=0).learn(envs, ...)}) and loaded with
    \code{MyAlgorithm.load(path)}.
  \item Add presets as a table \code{\{name: \{env\_id: preset\}\}} in a
    module of \pkg{marl\_algorithms.presets} and merge it into \code{PRESETS} in
    \code{presets/\_\_init\_\_.py}, where \code{get\_preset}, \code{compare} and
    \code{marl-train presets} look them up.
  \item Add tests in \code{tests/algorithms/}: action shapes on several
    environments, \code{policy()} with \code{env\_lib.evaluate}, the
    \code{TypeError} for an unsupported action kind, a short \code{learn} with
    shared and separate parameters, seed determinism, a \code{save}/\code{load}
    round trip and a loose learning check on a small task.
\end{enumerate}

A minimal skeleton, an independent advantage actor-critic for discrete
actions (one gradient step per rollout on the agents' own rewards):

% norun
\begin{pycode}
from dataclasses import dataclass
from typing import ClassVar

import torch
from torch import nn

from marl_algorithms.core.base import OnPolicyAlgorithm
from marl_algorithms.core.buffers import compute_gae
from marl_algorithms.core.config import OnPolicyConfig
from marl_algorithms.core.networks import CategoricalPolicy, PerAgent, mlp


@dataclass
class A2CConfig(OnPolicyConfig):
    """Options of IA2C (this sketch ignores the normalisation options)."""

    ent_coef: float = 0.01


class IA2C(OnPolicyAlgorithm):
    """Independent advantage actor-critic (discrete actions)."""

    name: ClassVar[str] = "ia2c"
    action_kinds: ClassVar[tuple[str, ...]] = ("discrete",)
    config_class: ClassVar[type[A2CConfig]] = A2CConfig

    def _build(self):
        cfg, spec = self.config, self.spec
        self.actor = PerAgent(
            lambda: CategoricalPolicy(self.agent_input_dim, spec.n_actions,
                                      cfg.hidden_sizes, cfg.activation),
            spec.n_agents, cfg.share_parameters).to(self.device)
        self.critic = PerAgent(
            lambda: mlp(self.agent_input_dim, 1, cfg.hidden_sizes,
                        cfg.activation),
            spec.n_agents, cfg.share_parameters).to(self.device)
        self.params = [*self.actor.parameters(), *self.critic.parameters()]
        self.optimizer = torch.optim.Adam(self.params, lr=cfg.lr)

    def _modules(self):
        return {"actor": self.actor, "critic": self.critic,
                "optimizer": self.optimizer}

    def act(self, obs, deterministic=False):
        inputs = self.agent_inputs(self.tensor(obs))
        with torch.inference_mode():
            if deterministic:
                return self.actor(inputs).argmax(-1).cpu().numpy()
            actions, _ = self.actor.call("sample", inputs, self.torch_rng)
        return actions.cpu().numpy()

    def _rollout_step(self, obs):
        return self.act(obs), {}                 # no extra rollout fields

    def _update(self, rollout):
        cfg = self.config
        inputs = self.agent_inputs(self.tensor(rollout["obs"]))
        values = self.critic(inputs).squeeze(-1)                  # (T, B, n)
        with torch.no_grad():
            next_inputs = self.agent_inputs(self.tensor(rollout["next_obs"]))
            next_values = self.critic(next_inputs).squeeze(-1)
        advantages, returns = compute_gae(
            rollout["agent_rewards"], values.detach().cpu().numpy(),
            next_values.cpu().numpy(), rollout["terminated"],
            rollout["truncated"], cfg.gamma, cfg.gae_lambda)
        actions = self.tensor(rollout["actions"], torch.int64)
        log_prob, entropy = self.actor.call("log_prob", inputs, actions)
        loss = (-(log_prob * self.tensor(advantages)).mean()
                - cfg.ent_coef * entropy.mean()
                + 0.5 * (values - self.tensor(returns)).pow(2).mean())
        self.optimizer.zero_grad(set_to_none=True)
        loss.backward()
        nn.utils.clip_grad_norm_(self.params, cfg.max_grad_norm)
        self.optimizer.step()
        return {"loss": loss.item(), "entropy": entropy.mean().item()}


# Registration (in marl_algorithms/registry.py) and a preset (in a module of
# marl_algorithms/presets/, merged into PRESETS):
# AlgorithmInfo("ia2c", "my_package.ia2c:IA2C", "on-policy", ("discrete",),
#               "independent advantage actor-critic", "<reference>")
PRESETS = {"ia2c": {"LineMsg-v0": {"num_envs": 32, "total_steps": 100_000,
                                   "config": {"lr": 1e-3}, "env_kwargs": {}}}}
\end{pycode}

% ---------------------------------------------------------------------------
\section{References}\label{sec:marl-references}

The source code of \pkg{marl\_algorithms} cites the following publications.

\begin{itemize}[leftmargin=1.5em,itemindent=-1.5em,label={}]
  \item J. Ackermann, V. Gabler, T. Osa and M. Sugiyama (2019).
    \emph{Reducing overestimation bias in multi-agent domains using double
    centralized critics.} NeurIPS Deep RL Workshop, arXiv:1910.01465.
  \item M. Andrychowicz et al.\ (2021). \emph{What matters in on-policy
    reinforcement learning? A large-scale empirical study.} ICLR.
  \item C. S. de Witt, T. Gupta, D. Makoviichuk, V. Makoviychuk,
    P. H. S. Torr, M. Sun and S. Whiteson (2020). \emph{Is independent
    learning all you need in the StarCraft multi-agent challenge?}
    arXiv:2011.09533.
  \item S. Fujimoto, H. van Hoof and D. Meger (2018). \emph{Addressing
    function approximation error in actor-critic methods.} ICML.
  \item T. P. Lillicrap et al.\ (2016). \emph{Continuous control with deep
    reinforcement learning.} ICLR.
  \item R. Lowe, Y. Wu, A. Tamar, J. Harb, P. Abbeel and I. Mordatch (2017).
    \emph{Multi-agent actor-critic for mixed cooperative-competitive
    environments.} NeurIPS.
  \item V. Mnih et al.\ (2015). \emph{Human-level control through deep
    reinforcement learning.} Nature 518, 529--533.
  \item T. Rashid et al.\ (2018). \emph{QMIX: monotonic value function
    factorisation for deep multi-agent reinforcement learning.} ICML.
  \item J. Schulman, P. Moritz, S. Levine, M. Jordan and P. Abbeel (2016).
    \emph{High-dimensional continuous control using generalized advantage
    estimation.} ICLR.
  \item J. Schulman, F. Wolski, P. Dhariwal, A. Radford and O. Klimov (2017).
    \emph{Proximal policy optimization algorithms.} arXiv:1707.06347.
  \item D. Silver et al.\ (2014). \emph{Deterministic policy gradient
    algorithms.} ICML.
  \item P. Sunehag et al.\ (2018). \emph{Value-decomposition networks for
    cooperative multi-agent learning based on team reward.} AAMAS.
  \item M. Tan (1993). \emph{Multi-agent reinforcement learning: independent
    vs.\ cooperative agents.} ICML.
  \item H. van Hasselt, A. Guez and D. Silver (2016). \emph{Deep
    reinforcement learning with double Q-learning.} AAAI.
  \item C. Yu, A. Velu, E. Vinitsky, J. Gao, Y. Wang, A. Bayen and Y. Wu
    (2022). \emph{The surprising effectiveness of PPO in cooperative
    multi-agent games.} NeurIPS Datasets and Benchmarks.
\end{itemize}

\endgroup
